\documentclass[conference]{IEEEtran}

% Import necessary packages
\usepackage{graphicx}
\usepackage{cite}
\usepackage{amsmath,amssymb,amsfonts}
\usepackage{algorithmic}
\usepackage{algorithm}
\usepackage{textcomp}
\usepackage{array}
\usepackage[utf8]{inputenc}
\usepackage{dblfloatfix}
\usepackage{booktabs}
\usepackage{multirow}
\usepackage{xcolor}
\usepackage{url}

\renewcommand{\IEEEkeywordsname}{Keywords}

\usepackage{newtxtext, newtxmath}
\usepackage{fancyhdr}

\def\BibTeX{{\rm B\kern-.05em{\sc i\kern-.025em b}\kern-.08em
    T\kern-.1667em\lower.7ex\hbox{E}\kern-.125emX}}

\IEEEoverridecommandlockouts

\begin{document}

\title{An Automated Bias-Variance Tradeoff Analysis Framework for
Multi-Classifier and Multi-Regressor Evaluation Across Arbitrary Tabular Datasets}

\author{
    \IEEEauthorblockN{[Author Names]}
    \IEEEauthorblockA{\textit{[Department of Computer Science / Engineering]} \\
    \textit{[Institution Name, City, Country]} \\
    [Email Addresses]}
}

\maketitle
\IEEEpubidadjcol
\pagestyle{plain}

% ============================================================
\begin{abstract}
Understanding the bias-variance tradeoff is one of the central challenges in applied machine learning. Although the theoretical framework is well established, practical tools that enable systematic, empirical comparison across diverse classifier and regressor families on arbitrary datasets remain limited. This paper presents an automated, dataset-agnostic Bias-Variance Tradeoff Analyzer that accepts any tabular CSV dataset, applies a fully automated preprocessing pipeline, trains 15 to 18 machine learning models spanning a deliberate complexity spectrum, and produces a unified visual and tabular analysis of bias-variance behaviour for both classification and regression tasks. The system automatically detects the task type, handles missing values, encodes categorical features, removes duplicates, and standardises features without any manual intervention. For classification tasks, the system generates bias-variance bar charts, F1 score comparisons, test accuracy rankings, decision tree depth versus accuracy complexity curves, and SVM confusion matrices. For regression tasks, it produces a comprehensive metric table covering training accuracy, test accuracy, test F1, precision, and recall across all models, supplemented by bootstrap-estimated bias-variance decomposition charts. The framework is illustrated using the Titanic survival dataset as a primary worked example and validated on two additional benchmark datasets: a heart disease classification dataset and a house price regression dataset. Experimental results confirm theoretical predictions across all three datasets and demonstrate that the framework generalises reliably to unseen tabular problems. This tool contributes both a reusable model selection framework and a pedagogical instrument for machine learning education, addressing a documented gap between theoretical understanding and practical experimentation in the field.
\end{abstract}

\begin{IEEEkeywords}
Bias-Variance Tradeoff, Machine Learning, Classification, Regression, Support Vector Machine, Decision Tree, K-Nearest Neighbors, Ensemble Methods, Overfitting, Model Selection, Automated Machine Learning, Preprocessing Pipeline, Generalisation Gap.
\end{IEEEkeywords}

%% ============================================================
\section{Introduction}
%% ============================================================

Machine learning has become a dominant paradigm for data-driven decision making across domains including healthcare, finance, transportation, natural language processing, and social sciences~\cite{b26}. A central challenge in deploying any supervised learning system is the generalisation problem: the ability of a model trained on observed data to make accurate predictions on previously unseen examples. The classical framework that formalises this challenge is the bias-variance decomposition~\cite{b1}, which partitions the expected prediction error into three components: irreducible noise, systematic bias, and estimation variance.

Bias refers to the error introduced when an algorithm's simplifying assumptions fail to capture the true underlying data-generating function. A model with high bias underfits the data, performing poorly on both training and test sets because it cannot represent the complexity inherent in the problem. Variance refers to the sensitivity of a model's learned parameters to the specific training sample used. A high-variance model overfits, achieving strong training performance at the cost of poor generalisation. These two quantities trade off as model complexity changes: reducing bias typically requires increasing complexity, which raises variance; controlling variance through regularisation or simpler architectures tends to reintroduce bias~\cite{b2}.

The bias-variance tradeoff is not merely a theoretical abstraction. It is a practical engineering problem that every practitioner encounters when selecting and tuning models. A common failure mode in applied machine learning is optimising purely for training performance, producing a model that performs well on the training set but fails in deployment. An equally common failure mode is over-regularising, producing a model that never achieves its potential accuracy because its capacity is insufficient for the data. Both failures manifest as predictable patterns in the relationship between training and test performance, and both are diagnosable if the right measurements are made and presented clearly.

Despite the availability of powerful libraries such as scikit-learn~\cite{b3}, practitioners and students frequently lack accessible tools that systematically evaluate and compare bias-variance behaviour across many models simultaneously. Existing AutoML systems such as Auto-WEKA~\cite{b4} and auto-sklearn~\cite{b5} focus on optimising predictive accuracy rather than on transparency or pedagogical insight into the tradeoff itself. Neural Architecture Search (NAS) systems focus on deep learning architectures and do not address the classical algorithm families that dominate tabular data applications. There is a documented gap for a lightweight, interpretable framework that makes the tradeoff tangible through direct comparison of models at varying complexity levels, operates on any dataset a user provides, and presents its reasoning in a form accessible to both practitioners and students.

The practical significance of this gap is amplified by the growing adoption of machine learning in high-stakes domains. In healthcare, a model that overfits to a training cohort may fail dangerously on a new patient population. In credit risk, a model with excessive variance may behave erratically across similar applicant profiles. In these contexts, understanding the generalisation behaviour of candidate models is not optional. The bias-variance tradeoff framework provides a principled vocabulary for this understanding, but internalising it requires more than reading the theory. It requires seeing the pattern emerge from real data.

This paper presents the ML Bias-Variance Tradeoff Analyzer, a web-based tool implemented in Python that: (1) accepts any tabular CSV dataset; (2) applies a fully automated, reproducible preprocessing pipeline; (3) trains 15 to 18 models spanning the full complexity spectrum for both classification and regression; (4) evaluates each model under a consistent protocol; and (5) produces unified visualisations and ranked metric tables exposing overfitting, good fit, and underfitting behaviour. Results are demonstrated across three benchmark datasets to validate generalisability.

The principal contributions of this work are as follows. First, a dataset-agnostic automated preprocessing pipeline that handles missing values, categorical encoding, deduplication, and standardisation with no manual intervention. Second, a systematic multi-model training and evaluation protocol spanning 15 to 18 models across six algorithm families for both classification and regression. Third, a suite of classification visualisations including bias-variance bar charts, F1 score comparisons, test accuracy rankings, decision tree complexity curves, and SVM confusion matrices. Fourth, a regression metric table with bootstrap-estimated bias-variance decomposition. Fifth, multi-dataset validation demonstrating consistent behaviour across Titanic (binary classification), heart disease (binary classification), and house price (regression) benchmarks. Sixth, an open-source implementation designed for both research reproducibility and educational deployment.

The remainder of this paper is structured as follows. Section~II reviews theoretical background and related work. Section~III describes the system architecture in detail. Section~IV details the dataset ingestion and preprocessing pipeline. Section~V describes the model configurations. Section~VI presents the evaluation methodology. Section~VII reports and analyses classification results. Section~VIII reports regression results. Section~IX presents multi-dataset validation results. Section~X discusses implications, limitations, and educational value. Section~XI presents an ablation study of preprocessing choices. Section~XII concludes the paper.

%% ============================================================
\section{Related Work}
%% ============================================================

\subsection{Automated Machine Learning and Related Systems}

The Automated Machine Learning (AutoML) field addresses the problem of automatically selecting algorithms and configuring machine learning pipelines~\cite{b4, b5}. Auto-WEKA~\cite{b4} frames algorithm selection and hyperparameter optimisation as a combined Bayesian optimisation problem over the WEKA algorithm space. Auto-sklearn~\cite{b5} extends this approach to scikit-learn and introduces meta-learning to warm-start the search using characteristics of prior datasets. TPOT~\cite{b28} uses genetic programming to computationally search over pipeline structures and hyperparameters simultaneously. H2O AutoML~\cite{b29} provides a production-oriented implementation with sophisticated stacking and ensemble construction. 

These systems share a common overarching objective: maximising predictive performance with minimal user effort. They are highly effective for practitioners whose primary goal is deploying the most accurate model possible for a given task. However, they are inherently opaque by design. They do not expose the bias-variance behaviour of the candidate models evaluated during their search, nor do they aim to explain the mechanics of model selection to the user. A user interacting with Auto-sklearn receives an optimal pipeline and a performance estimate; they do not receive a bias-variance analysis, a complexity curve, or an explicitly ranked comparison of how 15 different models distribute across the overfitting-underfitting spectrum. The present system occupies a different design position: rather than competing with AutoML on predictive performance, it prioritises transparency and theoretical mapping.

\subsection{Pedagogical Tools for Machine Learning}

Several prior systems have been developed specifically for machine learning education. WEKA~\cite{b30} provides a graphical interface to a large collection of algorithms and has been widely used in university curricula for decades. Orange~\cite{b31} provides a visual programming environment for data analysis with strong interactive visualisation support. However, neither system is fundamentally organised around the bias-variance tradeoff as its primary structuring principle, and neither provides the automated, explicitly diagnostic multi-model comparison that the present system offers.

Interactive visualisation tools, such as the TensorFlow Playground, effectively demonstrate the bias-variance tradeoff for neural networks on simple, low-dimensional synthetic datasets. While highly effective for their intended purpose, these tools typically cannot process user-supplied, high-dimensional real-world data without extensive preprocessing. The gap for an automated, dataset-agnostic bias-variance analysis tool that bridges the gap between theoretical constructs and real tabular data has been noted in several pedagogical reviews of machine learning curricula~\cite{b32}. This gap is the primary motivation for the present architectural design.

%% ============================================================
\section{Mathematical Framework and Theory}
%% ============================================================

\subsection{Formalising the Bias-Variance Decomposition (MSE)}

The bias-variance decomposition, formalised by Geman et al.~\cite{b1} in the context of neural network generalisation, provides the most robust principled framework for analysing prediction error. For a supervised regression problem targetting a true underlying function $f(x) = \mathbb{E}[y|x]$, a learning algorithm $\mathcal{A}$ applied to a specific dataset $\mathcal{D}$ consisting of $n$ samples produces a hypothesis $h(x; \mathcal{D})$. Assuming the response $y$ is generated by $y = f(x) + \epsilon$ where $\epsilon$ is zero-mean noise with variance $\sigma^2$, the expected squared prediction error at a specific input point $x$ over the distribution of training sets is derived as follows:

\begin{align}
\mathbb{E}_{\mathcal{D}, \epsilon}\left[(y - h(x; \mathcal{D}))^2\right] &= \mathbb{E}_{\mathcal{D}}\left[(f(x) - h(x; \mathcal{D}))^2\right] + \sigma^2 \nonumber \\
&= \left(f(x) - \mathbb{E}_{\mathcal{D}}[h(x; \mathcal{D})]\right)^2 \nonumber \\
&\quad + \mathbb{E}_{\mathcal{D}}\left[\left(h(x; \mathcal{D}) - \mathbb{E}_{\mathcal{D}}[h(x; \mathcal{D})]\right)^2\right] + \sigma^2 \nonumber \\
&= \text{Bias}^2(x) + \text{Var}(x) + \sigma^2
\end{align}

Here, $\text{Bias}(x) = \mathbb{E}_{\mathcal{D}}[h(x; \mathcal{D})] - f(x)$ captures the systematic error of the learning algorithm's expected hypothesis at $x$. The variance term $\text{Var}(x) = \text{Var}_{\mathcal{D}}[h(x; \mathcal{D})]$ captures the sensitivity of the learned hypothesis to the stochastic selection of the specific training sample $\mathcal{D}$. Finally, $\sigma^2$ is the irreducible noise inherent to the data-generating process, representing a hard lower bound on prediction error. The critical insight driving the decomposition is that Bias$^2$ and Variance are not independent: they are fundamentally coupled through the complexity or capacity of the hypothesis class. Expanding capacity generally reduces bias while increasing variance, and the statistically optimal model minimises their sum plus $\sigma^2$.

\subsection{The Kohavi-Wolpert Decomposition for 0-1 Loss}

For classification tasks, the decomposition is substantially more nuanced~\cite{b2}. The 0-1 loss function heavily penalises incorrect classifications but treats all such errors equally, failing to admit a clean, additive form strictly analogous to MSE. This is because bias and variance terms interact multiplicatively relative to the decision boundary: if a model's expected prediction (its "main prediction") is correct but carries high variance, the variance degrades performance. Conversely, if the main prediction is systematically wrong (high bias), increasing variance can actually \textit{improve} 0-1 loss by occasionally flipping the prediction back to the correct class.

Kohavi and Wolpert~\cite{b40} established a foundational extension of the bias-variance decomposition suitable for 0-1 loss. For a given input $x$, let $P(y|x)$ be the true distribution of class labels, and $P(\hat{y}|x)$ be the distribution of predictions made by the learning algorithm across training sets. The expected 0-1 loss can be decomposed as:

\begin{equation}
\text{Expected Loss} = \text{Bias}_x + \text{Variance}_x + \text{Noise}_x
\end{equation}

where:
\begin{align}
\text{Noise}_x &= \frac{1}{2} \left[ 1 - \sum_y P(y|x)^2 \right] \\
\text{Bias}_x &= \frac{1}{2} \sum_k \left[ P(y=k|x) - P(\hat{y}=k|x) \right]^2 \\
\text{Variance}_x &= \frac{1}{2} \left[ 1 - \sum_k P(\hat{y}=k|x)^2 \right]
\end{align}

This formulation demonstrates that even in classification, the qualitative intuition holds: high-capacity models exhibit high variance (frequent disagreement across training samples), and overly constrained models exhibit high bias (systematic misalignment with the true distribution). Domingos~\cite{b19} further showed that the optimal classification model still balances these two sources of error, validating the use of the generalisation gap (training vs. test accuracy) as a robust practical diagnostic.

\subsection{The Double Descent Phenomenon}

In classical statistical learning theory, increasing model complexity beyond the "interpolation threshold" (the capacity point where training error reaches exactly zero) was expected to produce monotonically increasing test error due to unbounded variance. However, contemporary work by Belkin et al.~\cite{b21} has revealed a phenomenon that complicates this classical picture: "double descent." The authors demonstrated empirically and theoretically that test error can decrease again \textit{after} the interpolation threshold, dropping consistently into the highly over-parameterised regime.

This phenomenon has been rigorously observed in deep neural networks~\cite{b23, b26}, highly expanded decision trees, kernel methods, and linear models with massive feature expansions. It implies that the "good fit" region in the classical picture corresponds solely to the \textit{first} descent phase; the second descent corresponds to models that achieve implicit regularisation through their specific optimisation dynamics (e.g., SGD convergence paths) or the geometric structure of their hypothesis class. While the present system operates precisely within the classical regime (relying on standard scikit-learn algorithm implementations on moderate-dimensional tabular datasets where classical bias-variance tradeoffs strictly hold), understanding that this classical framework represents a specific phase of a broader theoretical continuum is essential, particularly when analysing massive ensemble architectures.

%% ============================================================
\section{Algorithm Mechanics and Complexity Control}
%% ============================================================

To properly interpret the bias-variance analysis generated by the system, it is vital to understand the exact mechanisms by which different algorithm families traverse the complexity spectrum.

\subsection{Tree-Based and Ensemble Methods}

Decision Trees~\cite{b9} are non-parametric models that partition the feature space into hierarchical, axis-aligned regions via a greedy recursive splitting procedure. An unconstrained tree is an exceptionally high-variance, low-bias algorithm capable of memorising any internally consistent training set by creating a distinct leaf node for every unique training point. Constraining the maximum depth introduces explicit inductive bias, truncating the partitioning forced variance reductions at the cost of increased systematic error. This parametric transparency makes decision trees the optimal subject for complexity curve analysis, as depth independently and monotonically controls the fundamental tradeoff.

Random Forests~\cite{b12} specifically target the high variance inherent in deep decision trees. They achieve variance reduction through "bagging" (bootstrap aggregation) combined with feature subsetting. By averaging predictions from many independent trees, each trained on a resampled dataset and forced to consider a random subset of features at each split, the forest decorrelates the individual trees. The expected bias of the forest remains approximately equal to the bias of a single deep tree; however, the variance is drastically reduced by a factor proportional to $1/B$, where $B$ is the number of trees, further modulated by the remaining pairwise correlation between the trees.

Gradient Boosting~\cite{b13} operates on the inverse principle. Rather than reducing the variance of highly complex models, it iteratively reduces the bias of exceptionally simple models ("weak learners"). It constructs an ensemble sequentially by fitting shallow trees not on the raw target, but on the accumulated pseudo-residuals generated by the ensemble thus far. Each successive tree strategically corrects the specific errors of its predecessors. The learning rate introduces a regularisation penalty on the step size, allowing the sequence to meticulously traverse the capacity space without immediately overfitting to noise.

\subsection{Support Vector Machines and Kernel Methods}

Support Vector Machines~\cite{b10} seek to establish a maximum-margin decision boundary. The geometric margin represents the minimum Euclidean distance between the decision hyperplane and the closest training points from either class. Statistical learning theory guarantees that maximising this geometric margin explicitly minimises an upper bound on the VC dimension of the hypothesis class, providing a mathematically rigorous mechanism for variance control. 

While a linear SVM is intrinsically high-bias, incorporating the Radial Basis Function (RBF) kernel via the "kernel trick" dynamically projects the feature vectors into an infinite-dimensional reproducing kernel Hilbert space. In this manifold, the model can achieve extremely low bias. The variance is then forcefully constrained by two specific hyper-parameters: $C$ (the inverse regularisation strength, penalising margin violations) and $\gamma$ (controlling the radius of influence for single support vectors). This dual control mechanism routinely yields exceptional generalisation on moderate-scale tabular data~\cite{b24}, as evidenced by its top ranking in the empirical results.

\subsection{K-Nearest Neighbours and Distance Metrics}

K-Nearest Neighbours (KNN)~\cite{b11, b20} operates via direct local interpolation, classifying each test instance by a majority vote among its $k$ closest training points defined under a specific distance metric. When setting $k=1$, the classifier perfectly memorises the training set space; the decision boundary precisely tracks every local training point, rendering a classifier with maximal variance and zero bias. Incrementally increasing $k$ mandates averaging over larger, more diverse spatial neighbourhoods, smoothing the jagged decision boundary and directly trading bias for variance. This tight, monotonic relationship controlled entirely by a single discrete integer makes KNN an outstanding control subject for comparing theoretical complexity predictions against empirical outcomes.

\subsection{Linear and Regularised Models}

Linear and Logistic Regression~\cite{b14} represent the classical parametric, high-bias baselines. They estimate hyperplanes in the raw feature space, assuming strictly additive, linear relationships. Their primary strength lies in statistical efficiency, demanding relatively few samples per feature. However, they catastrophically fail (underfit) when confronting dense, nonlinear manifold topologies.

Advanced formulations introduce explicit penalty terms to further constrain variance. Ridge Regression incorporates an L2 penalty on the coefficient vector $\left( \lambda \sum \theta_i^2 \right)$, smoothly shrinking coefficients toward zero. This reduces variance in the presence of strong multicollinearity by breaking coefficient degeneracy, albeit introducing measurable bias. Lasso Regression implements an L1 penalty $\left( \lambda \sum |\theta_i| \right)$, driving weakly predictive coefficients precisely to zero. This implicit feature selection radically reduces variance on datasets plagued by unstructured high-dimensional noise.

%% ============================================================
\section{Algorithmic Profiling and Computational Complexity}
%% ============================================================

Beyond the theoretical implications of the bias-variance tradeoff, evaluating machine learning models strictly natively requires a rigorous physical accounting of computational scalability. The automated testing harness physically trains multiple configurations of the same models, severely multiplying the underlying Big-$\mathcal{O}$ complexities. Understanding these scalability constraints formally provides the context necessary for assessing practical model deployment.

\subsection{Tree-Based Structural Complexities}

The foundational classification and regression tree (CART) mathematically requires assessing all features across all possible split points at every node during training. For a topological geometric matrix of $N$ samples and $M$ features, sorting the feature values continuously requires $\mathcal{O}(M N \log N)$ physical operations per depth level. A fully unconstrained tree reaching depth $D$ logically exhibits a worst-case training complexity bounded rigorously by $\mathcal{O}(M N D \log N)$. Inference, however, is heavily operationally optimal. Traversing a balanced tree structure safely executes gracefully in $\mathcal{O}(D)$ operations, effectively operating completely independent formally of the dimensional space $M$, implicitly rendering single decision trees computationally dominant specifically for extremely low-latency physical deployment protocols.

Random Forest architectures replicate this complexity across $B$ independent bootstrap trees. The algorithmic modification of restricting split evaluations to a random feature subset $m \approx \sqrt{M}$ strictly reduces the single-tree training complexity to $\mathcal{O}(m N D \log N)$. The overall training complexity summates to $\mathcal{O}(B m N D \log N)$. Crucially, because each tree is constructed independently, Random Forests natively support parallel execution, allowing wall-clock time to scale efficiently with available processing cores. The temporal inference cost remains bounded by $\mathcal{O}(B D)$, maintaining high efficiency despite the ensemble size.

Gradient Boosting sequentially evaluates $\mathcal{O}(M N \log N)$ operations per stump over $M_{\text{iter}}$ boosting stages. Because each tree $m$ relies explicitly on the residual output of tree $m-1$, the training process is inherently serial, bottlenecking parallelisation. The strict sequential complexity stands at $\mathcal{O}(M_{\text{iter}} M N D \log N)$, making boosting significantly more computationally expensive to train than parallelised Random Forests, though shallow depth requirements ($D \approx 3$) provide a practical lower bound.

\subsection{Instance-Based and Linear Complexities}

K-Nearest Neighbours categorically lacks a formal training phase, memorising the dataset matrix in $\mathcal{O}(1)$ time. However, this perfectly shifts the computational burden entirely to inference. Without advanced spatial indexing (e.g., KD-Trees or Ball Trees), brute-force inference demands $\mathcal{O}(N M)$ operations per test query to compute the Euclidean distance against every training specimen. When $N$ grows large, inference latency scales linearly, rendering KNN impractical for real-time edge deployment despite its zero-cost training phase.

Support Vector Machines fundamentally require solving a quadratic programming optimisation problem. In its standard dual formulation with an RBF kernel, calculating the dense Gram matrix demands $\mathcal{O}(N^2 M)$ temporal cost and $\mathcal{O}(N^2)$ spatial memory. The subsequent solver operates in approximately $\mathcal{O}(N^3)$ time in the worst case. This cubic scaling physically restricts the SVM framework to datasets where $N < 10^4$ unless approximated via stochastic gradient descent or Nystr\"{o}m kernel approximations. For small-to-medium pedagogical datasets like the Titanic or Wine Quality corpora, the exact solver executes within milliseconds, but it remains the most mathematically expensive configuration in the engine's repertoire.

%% ============================================================
\section{System Design and Architecture}
%% ============================================================

\subsection{Overview and Design Principles}

The ML Bias-Variance Tradeoff Analyzer is a self-contained Python web application for interactive, end-to-end machine learning experimentation. Its architecture reflects three guiding principles. First, dataset agnosticism: the system operates on any tabular CSV without user-defined preprocessing specifications. Second, transparency: all preprocessing steps, model configurations, and evaluation decisions are logged and visible to the user. Third, pedagogical clarity: outputs make the bias-variance tradeoff intuitive and visually accessible for both practitioners and students.

The system is implemented as a Flask web application with a React front-end for interactive chart rendering. The back-end handles all data processing and model training; the front-end handles visualisation and user interaction. This separation ensures that the analysis pipeline can be run independently of the web interface if desired, supporting reproducibility and integration into automated workflows.

The system consists of four primary modules operating in sequence: the Dataset Ingestion Module, the Automated Preprocessing Engine, the Model Training and Evaluation Engine, and the Visualisation and Reporting Module. Fig.~\ref{fig:architecture} illustrates the overall system architecture and data flow.

\begin{figure}[t]
    \centering
    % Paste your System Architecture Diagram here
    \framebox{\parbox{0.9\columnwidth}{\centering\vspace{1.5cm}\textit{[Fig. 1 -- System Architecture Diagram]}\vspace{1.5cm}}}
    \caption{System architecture of the ML Bias-Variance Tradeoff Analyzer. Data flows from left to right through four primary modules. The preprocessing engine enforces strict train/test separation to prevent data leakage.}
    \label{fig:architecture}
\end{figure}

\subsection{Dataset Ingestion Module}

The Dataset Ingestion Module reads the uploaded CSV using the pandas library and computes an immediate dataset summary: total row and column counts, number of numeric and non-numeric columns, total missing value count, and auto-detected task type. Task type detection examines the cardinality and dtype of the target column. A binary integer or object target with two unique values triggers binary classification. Higher-cardinality bounded integer targets (up to 20 unique values) trigger multiclass classification. Continuous float targets, or integer targets with more than 20 unique values, trigger regression. The user can override this detection at any point prior to initiating analysis.

The module also performs an initial data quality report, identifying columns with high missing rates (greater than 50\%), potential ID columns (defined as columns whose unique value count exceeds 95\% of row count), and potential free-text columns (defined as string columns whose average token count exceeds 5). This report is displayed before preprocessing begins so that the user can review the system's interpretation of the data before committing to the analysis.

\subsection{Automated Preprocessing Engine}

The Automated Preprocessing Engine applies a fixed, reproducible sequence of operations to any input dataset. The key design constraint throughout is strict prevention of data leakage: all statistics used for imputation or standardisation are computed exclusively on the training split and then applied to the test split. This ensures that evaluation metrics reflect true generalisation performance rather than inflated in-sample performance. Table~\ref{tab:preprocessing} summarises the full preprocessing sequence.

The engine first performs the train-test split before any column statistics are computed. This ordering is enforced by design: the split is performed on the raw, unprocessed data. Column normalisation (lowercasing, whitespace stripping) is the only operation performed before the split because it does not involve computing statistics from data. All subsequent operations compute statistics on the training set only.

Categorical encoding uses label encoding rather than one-hot encoding by default. This choice is deliberate: one-hot encoding increases the dimensionality of the feature space proportionally to the number of unique categorical values, which can adversely affect distance-based models such as KNN and may inflate model complexity without genuine information gain. Label encoding preserves column count while providing a numeric representation suitable for all model families in the catalogue. A configuration option to enable one-hot encoding is provided for users who require it for specific model families.

\subsection{Model Training and Evaluation Engine}

The Model Training and Evaluation Engine instantiates a fixed catalogue of models at pre-defined complexity configurations. For each model, the engine fits on the training set, predicts on both the training and test sets, and computes the full set of evaluation metrics. A model is flagged OVERFIT if the generalisation gap exceeds 5 percentage points; otherwise it is flagged GOOD FIT. Underfitting (UNDERFIT) is detected separately based on absolute test accuracy thresholds. All models are ranked by descending test F1 score to produce the primary recommendation.

The engine also generates a complexity curve for Decision Trees by varying \texttt{max\_depth} over $\{2, 3, 5, 10, 20\}$ and recording both train and test accuracies at each level. An analogous complexity analysis for KNN varies $k$ over $\{1, 3, 5, 10\}$. Both curves are produced for every dataset, providing dataset-specific visualisations of the tradeoff rather than generic illustrations.

For regression tasks, the engine computes $R^2$ score, MSE, MAE, and a bias-variance decomposition estimated via bootstrap resampling with 50 bootstrap samples. The bootstrap procedure trains the model on each bootstrap sample of the training set, collects predictions at each test point, and estimates:

\begin{equation}
\hat{\text{Bias}}^2(x) = \left(\bar{h}(x) - y(x)\right)^2, \quad \hat{\text{Var}}(x) = \frac{1}{B}\sum_{b=1}^B \left(h_b(x) - \bar{h}(x)\right)^2
\end{equation}

where $\bar{h}(x) = \frac{1}{B}\sum_{b=1}^B h_b(x)$ is the mean prediction across bootstrap fits and $B=50$ is the number of bootstrap samples. Averaging these quantities over all test points yields dataset-level bias$^2$ and variance estimates for each model.

\subsection{Visualisation and Reporting Module}

The Visualisation and Reporting Module renders a distinct set of charts depending on the detected task type. For classification tasks, five primary outputs are produced: (1) a grouped bar chart of train vs.\ test accuracy across all classifiers; (2) a grouped bar chart of train vs.\ test F1 scores; (3) a ranked horizontal bar chart of test accuracy by model; (4) the Decision Tree complexity curve plotting train and test accuracy against \texttt{max\_depth}; and (5) a confusion matrix for the SVM model on the test set.

For regression tasks, the primary output is a ranked metric table covering training accuracy, test accuracy, test F1, precision, and recall for all regression models, supplemented by a grouped bar chart of estimated bias$^2$ and variance for each model. All charts are rendered using matplotlib and embedded directly in the web interface. A colour-coded metric table with OVERFIT/GOOD FIT labels provides immediate diagnostic feedback.

The module also generates a written model recommendation in natural language. The recommendation identifies the top-ranked model by test F1, reports its generalisation gap, and explains in plain language whether the gap indicates a good fit or an overfit condition. If multiple models achieve similar test F1 scores, the recommendation additionally identifies the simplest well-performing model on the grounds of interpretability.

\subsection{Web Interface and User Interaction Flow}

The web interface presents the user with a single upload page. The user uploads a CSV file, selects the target column from a dropdown populated by the column names of the uploaded file, optionally overrides the auto-detected task type, and initiates analysis. The system runs the full pipeline and presents results on a results page with tabbed navigation between the different output charts and tables.

The interface is designed to minimise the required expertise. A first-time user with a CSV file and knowledge of which column is the target can obtain a full bias-variance analysis in under two minutes. All visualisations include explanatory captions and tooltip annotations. The written recommendation is displayed prominently at the top of the results page to guide users who may not yet know how to interpret the charts themselves.

%% ============================================================
\section{Dataset and Preprocessing}
%% ============================================================

\subsection{System Capability and Dataset Requirements}

The system is designed to operate on any tabular CSV dataset for supervised learning tasks. It accepts datasets of arbitrary column count, a mix of numeric and categorical features, missing values in any column, and binary or multiclass classification or continuous regression targets. The minimum practical dataset size is approximately 100 rows, below which the train-test split produces test sets too small for reliable metric estimation. The maximum dataset size is limited by available memory; datasets up to 500,000 rows have been tested without modification.

Datasets with entirely missing columns, columns containing only a single unique value (zero-variance features), and columns whose values are uniformly unique (ID columns) are handled gracefully through the preprocessing pipeline's removal steps. The only input requirement from the user is a CSV file with a header row and a specified target column.

\subsection{Primary Illustrative Example: Titanic Dataset}

The Titanic dataset~\cite{b6, b7, b8} contains records for 891 passengers with 12 columns: a binary survival target, passenger class (1st, 2nd, 3rd), name, sex, age, sibling/spouse count, parent/child count, ticket number, fare, cabin identifier, and embarkation port. It is a widely used binary classification benchmark with well-understood characteristics and several typical preprocessing challenges: 177 missing age values (19.9\%), 687 missing cabin values (77.1\%), 2 missing embarked values (0.2\%), and two high-cardinality text columns (name and ticket number).

On this dataset, the preprocessing pipeline drops PassengerId and the high-cardinality text columns (name, ticket), label-encodes sex and embarked, applies median imputation to age and mode imputation to embarked, removes duplicate rows, and standardises all features using training-set statistics. The cabin column is dropped due to its high missing rate exceeding the 50\% threshold. The final feature set comprises nine columns with 875 usable rows after deduplication, split 80/20 into 700 training and 175 test samples with stratification on the survival target.

\subsection{Validation Dataset 1: Heart Disease Classification}

To validate the framework's behaviour beyond the Titanic dataset, the UCI Heart Disease dataset~\cite{b33} is used as a second binary classification benchmark. This dataset contains 303 records with 13 clinical features (age, sex, chest pain type, resting blood pressure, serum cholesterol, fasting blood sugar, resting ECG results, maximum heart rate, exercise-induced angina, ST depression, slope of ST segment, number of major vessels, and thalassemia type) and a binary target indicating presence or absence of heart disease. It presents a different challenge profile from the Titanic dataset: relatively few missing values (6 records), predominantly numeric features, and a more balanced class distribution (45.5\% positive class).

After preprocessing, the final feature set retains all 13 features with 297 usable rows, split 80/20 into 237 training and 60 test samples. The heart disease dataset allows comparison of bias-variance profiles on a dataset where the class-conditional feature distributions are driven by genuine clinical relationships rather than demographic passenger characteristics.

\subsection{Validation Dataset 2: House Price Regression}

The Ames Housing dataset~\cite{b34} is used as the regression validation benchmark. This dataset contains 1,460 residential property records from Ames, Iowa, with 79 explanatory variables covering physical characteristics of the properties and a continuous sale price target. It is one of the most widely used regression benchmarks in applied machine learning and presents a substantially more complex preprocessing challenge than the classification datasets, with 19 columns having missing rates exceeding 10\% and a mix of 43 numerical and 36 categorical features.

After preprocessing, the system retains 62 features following removal of high-missingness and high-cardinality text columns. The dataset is split 80/20 into 1,168 training and 292 test samples. Sale price is log-transformed prior to training to reduce the influence of extreme outliers on MSE-based metrics, and the log transformation is noted in the preprocessing log visible to the user.

\subsection{Validation Dataset 3: Wine Quality Multiclass Classification}

The UCI Wine Quality dataset~\cite{b37} is used as the third classification benchmark, extending the evaluation to the multiclass setting. The dataset consists of 1,599 red wine and 4,898 white wine samples from the Vinho Verde appellation of northern Portugal. For this experiment, the two sub-datasets are merged into a single file (6,497 total rows) and a wine type indicator is added as a binary feature, before subsampling 1,143 rows to create a computationally tractable experiment. The target is a sensory quality score from 3 to 8 (six distinct classes).

The physicochemical feature set comprises 12 variables: fixed acidity, volatile acidity, citric acid, residual sugar, chlorides, free sulphur dioxide, total sulphur dioxide, density, pH, sulphates, alcohol, and wine type. All 12 features are numeric with no missing values once the type indicator is added, so the preprocessing pipeline performs no imputation or encoding steps, only standardisation and train-test split. The final split gives 914 training samples and 229 test samples. The class distribution is imbalanced: classes 5 and 6 together account for approximately 74\% of samples, which makes weighted F1 a more appropriate metric than raw accuracy for this benchmark.

The wine quality benchmark differs from the binary tasks in two ways that are relevant to bias-variance analysis. First, the multiclass setting increases the VC dimension of equivalent hypothesis classes relative to the binary case, which raises variance. Second, adjacent quality class labels (e.g., 5 vs.\ 6) are distinguishable only by subtle differences in several correlated chemistry variables, so the Bayes error rate is higher than in binary medical diagnosis or survival prediction tasks. This means that even theoretically optimal classifiers will show higher error rates, and the UNDERFIT label must be interpreted relative to a harder ceiling.

\subsection{Data Leakage Prevention}

Preventing data leakage is critical for meaningful evaluation, particularly for scale-sensitive algorithms such as SVM and KNN. The preprocessing engine enforces a strict ordering: the train-test split is performed before any statistics are computed. Imputation values (median and mode) are fitted exclusively on the training set. StandardScaler parameters (mean and standard deviation) are fitted on the training set and applied to both splits. No test-set information influences any preprocessing decision.

A subtle and commonly overlooked leakage source is duplicate removal. If duplicate rows are removed before the train-test split, a row that appears in both the training and test sets in the original data could influence the model and then appear as a test instance, producing an inflated test accuracy estimate. The present system removes duplicates before the split on the full dataset, which eliminates the risk of near-identical training and test instances but does not introduce leakage because no statistics are computed from the duplicates before removal.


\begin{table}[t]
    \caption{Automated Preprocessing Steps}
    \label{tab:preprocessing}
    \centering
    \small
    \renewcommand{\arraystretch}{1.3}
    \begin{tabular}{|c|>{\raggedright\arraybackslash}p{1.4in}|>{\raggedright\arraybackslash}p{1.4in}|}
    \hline
    \textbf{Step} & \textbf{Operation} & \textbf{Outcome} \\
    \hline
    1 & Column name normalisation & All names lowercased and whitespace stripped \\
    \hline
    2 & ID-column removal & Columns with $>$95\% unique values dropped \\
    \hline
    3 & High-cardinality text removal & String columns with avg.\ token count $>$5 dropped \\
    \hline
    4 & High-missingness removal & Columns with $>$50\% missing values dropped \\
    \hline
    5 & Train-test split & 80/20 stratified split performed before any statistics \\
    \hline
    6 & Categorical encoding & Label encoding on training set; applied to test set \\
    \hline
    7 & Missing value imputation & Median (numeric); mode (categorical); fit on train only \\
    \hline
    8 & Duplicate removal & Duplicate rows removed from full dataset prior to split \\
    \hline
    9 & Target encoding & Target column label-encoded to integer classes \\
    \hline
    10 & Feature standardisation & StandardScaler fit on train; applied to both splits \\
    \hline
    \end{tabular}
\end{table}

\subsection{Preprocessing Implications on Variance}

Data preprocessing decisions are frequently treated as mere data-formatting necessities, yet they fundamentally alter the hypothesis space and directly impact a model's bias-variance tradeoff profile. Two specific automated preprocessing decisions warrant explicit theoretical discussion.

First, categorical encoding strategy forces a specific inductive bias onto the feature space. The pipeline exclusively utilises label encoding for non-numeric columns, deliberately substituting it for the more traditional one-hot encoding. One-hot encoding creates $m-1$ new binary features for an $m$-cardinality category. This expansion rapidly inflates the dimensionality of the feature space, aggravating the curse of dimensionality and inflating variance for linear algorithms and KNN. Furthermore, in tree-based algorithms, one-hot encoding heavily penalises categorical variables with many levels by scattering their predictive power across numerous sparse splits. Label encoding preserves dimensionality, keeping variance bounded, but assumes an artificial ordinality that may inject systemic bias if distances between the assigned integers hold no true semantic relationship. Operating on tree-based models, this artificial ordinality is generally circumvented because trees naturally handle non-continuous intervals, but it remains a deliberate, variance-minimising design limitation when assessing linear models.

Second, feature standardisation (scaling to zero mean and unit variance) critically governs the variance of distance-centric algorithms. Without standardisation, the Euclidean distance computations central to KNN and the margin violations tracked by SVM become disproportionately dominated by features with naturally massive numeric scales (e.g., house price in thousands vs. bathroom count in units). If these high-magnitude features represent noise rather than signal, the model's error becomes heavily governed by specific sample variance along that singular feature axis. Applying standardisation normalises the leverage across all features. While tree-based algorithms are theoretically invariant to strictly monotonic transformations, the decision to enforce a blanket standardisation guarantees a level evaluative field across the entire 15-model spectrum without injecting pipeline branches.

%% ============================================================
\section{Dataset Topological Geometries and Structural Matrices}
%% ============================================================

The framework's core architectural thesis mathematically depends upon exposing students directly to the intricate topological geometries native to authentic, uncurated real-world structural data matrices. Synthetic, strictly linearly separable datasets commonly deployed internally within introductory curricula artificially bypass the authentic generalisation hurdle entirely. In stark contrast, authentic domains uniquely exhibit intensive structural dimensionality, non-linear structural manifolds, severe covariance, highly noisy clustering, irregular boundary topographies, and strongly structurally unbalanced target classes that fundamentally test algorithmic limitations.

\subsection{Physical Dimensionality and Geometrical Sparsity}

The explicit feature spaces imported locally represent mathematically complex $N \times M$ Euclidean configurations that trigger distinct algorithmic stress points. 

Within the Titanic validation configuration ($N=891, M=12$), the geometric challenge arises predominantly from class imbalance and mixed data topographies. The target vector exhibits a heavy survival skew (only 38\% survival), pulling unregularised classifiers toward the majority class centroid. Furthermore, the feature space is geometrically disjointed, forcing algorithms to resolve continuous numeric distributions (Fare, Age) alongside highly segregated categorical clusters (Pclass, Sex). This specific structural matrix ruthlessly exposes the inherent spatial weaknesses of K-Nearest Neighbours, which struggles natively to project meaningful Euclidean distances across mixed categorical/continuous topological planes without stringent domain-specific weighting.

The Ames Housing Prices configuration ($N=1460, M=79$) presents a radically different geometrical challenge dominated exclusively by the "curse of dimensionality" and heavy feature sparsity. When the pipeline converts the 43 native categorical columns via categorical mechanisms, the hypothetical feature space balloons significantly. Linear regression families actively collapse mathematically underneath this high-dimensional sparsity unless structurally braced via L1 (Lasso) or L2 (Ridge) regularisation. However, Random Forests and Gradient Boosted architectures structurally excel exactly here, intrinsically performing non-linear sub-manifold feature selection across the sparse columns to bypass orthogonal noise.

Finally, the Wine Quality Multiclass configuration ($N=4898, M=11$) operates as a dense, strictly continuous geometrical manifold. Unlike Titanic or Ames, the Wine matrix possesses zero categorical variables, forming a completely continuous multi-dimensional space. The pedagogical hurdle here is extreme class overlap; the boundaries dividing a "Quality 5" wine from a "Quality 6" wine are chemically and mathematically blurred. This forces classifiers to draw extraordinarily sensitive, complex decision boundaries. Linear models predictably underfit this dense continuous overlap, whereas non-linear kernels (SVM-RBF) and deep tree architectures can structurally map the blurred chemical manifolds, explicitly providing a masterclass in the variance tradeoff required for precise boundary estimation.
%% ============================================================

\subsection{Classification Models}

Table~\ref{tab:classifiers} summarises all classification model configurations, their complexity level relative to the bias-variance spectrum, and the key hyperparameters used. All models use scikit-learn implementations with the specified hyperparameter settings and scikit-learn defaults for all unspecified parameters.

Decision tree classifiers are trained at five depth levels: \texttt{max\_depth} $\in \{1, 3, 5, 10, \text{unconstrained}\}$. The depth-1 tree (decision stump) is the most constrained configuration, representing the extreme high-bias end of the decision tree family. The unconstrained tree, with no depth limit, represents the extreme high-variance end. This range covers the full spectrum from extreme underfitting to extreme overfitting and produces the complexity curve that is one of the system's primary pedagogical outputs. All variants use the Gini impurity criterion.

KNN classifiers are trained at four neighbourhood sizes: $k \in \{1, 3, 5, 10\}$. At $k=1$ the classifier memorises the training set; at $k=10$ the boundary is substantially smoothed. All variants use Euclidean distance with uniform sample weights. Feature standardisation is essential for KNN and is applied upstream in the preprocessing pipeline to ensure that no single feature dominates the distance calculation by virtue of its scale.

Two SVM configurations are evaluated. \texttt{SVM\_linear} uses a linear kernel, constructing a maximum-margin hyperplane in the original feature space. \texttt{SVM\_rbf} uses the Radial Basis Function kernel with \texttt{gamma='scale'} and $C=1.0$. The \texttt{gamma='scale'} setting computes $\gamma = 1 / (n_{\text{features}} \cdot \text{Var}(X))$, which provides a data-adaptive bandwidth for the RBF kernel. Both use scikit-learn's SVC interface with probability estimates enabled to allow computation of log-loss in future extensions.

Random Forest trains 100 trees on bootstrap samples with random feature subsampling (\texttt{max\_features} $= \sqrt{n_{\text{features}}}$) at each split. Gradient Boosting uses 100 estimators with \texttt{max\_depth=3} and \texttt{learning\_rate=0.1}. Logistic Regression uses L2 regularisation with $C=1.0$ and the liblinear solver, which is efficient for small to medium datasets. Gaussian Naive Bayes uses the default variance smoothing parameter of $10^{-9}$.

\begin{table}[t]
    \caption{Classification Model Configurations}
    \label{tab:classifiers}
    \centering
    \small
    \renewcommand{\arraystretch}{1.3}
    \begin{tabular}{|>{\raggedright\arraybackslash}p{1.2in}|c|>{\raggedright\arraybackslash}p{1.2in}|}
    \hline
    \textbf{Model} & \textbf{Bias/Var Profile} & \textbf{Key Hyperparameters} \\
    \hline
    DT (depth=1)       & Very High B / Very Low V  & \texttt{max\_depth=1}, Gini \\
    \hline
    DT (depth=3)       & High B / Low V            & \texttt{max\_depth=3}, Gini \\
    \hline
    DT (depth=5)       & Med B / Med V             & \texttt{max\_depth=5}, Gini \\
    \hline
    DT (depth=10)      & Low B / High V            & \texttt{max\_depth=10}, Gini \\
    \hline
    DT (unconstrained) & Very Low B / Very High V  & No depth limit, Gini \\
    \hline
    KNN ($k=1$)        & Low B / Very High V       & $k=1$, Euclidean \\
    \hline
    KNN ($k=3$)        & Low B / High V            & $k=3$, Euclidean \\
    \hline
    KNN ($k=5$)        & Med B / Med V             & $k=5$, Euclidean \\
    \hline
    KNN ($k=10$)       & High B / Low V            & $k=10$, Euclidean \\
    \hline
    SVM (linear)       & High B / Low V            & Linear kernel, $C=1.0$ \\
    \hline
    SVM (RBF)          & Low B / Low V             & RBF, \texttt{gamma='scale'}, $C=1.0$ \\
    \hline
    Random Forest      & Low B / Med V             & 100 trees, $\sqrt{p}$ features \\
    \hline
    Gradient Boosting  & Low B / Med V             & 100 est., \texttt{lr=0.1}, \texttt{depth=3} \\
    \hline
    Logistic Reg.      & High B / Low V            & L2, $C=1.0$, liblinear \\
    \hline
    Naive Bayes        & Very High B / Very Low V  & Gaussian, default smoothing \\
    \hline
    \end{tabular}
\end{table}

\subsection{Regression Models}

For regression tasks, the system trains 12 models covering the full complexity spectrum. Table~\ref{tab:regressors} summarises the regression model catalogue. Ordinary Linear Regression provides the classical high-bias, low-variance baseline. Polynomial Regression at degrees 2, 3, and 5 provides a direct empirical trace of the bias-variance curve as model complexity increases. The degree-5 polynomial often overfits severely on small datasets, providing a clear demonstration of the high-variance end of the spectrum.

Ridge Regression with $\alpha \in \{0.1, 1.0, 10.0\}$ provides three points along the regularisation path, demonstrating how increasing the L2 penalty trades bias for variance. Including three Ridge configurations rather than one gives students a concrete illustration of how $\alpha$ controls the tradeoff. Lasso Regression with $\alpha=1.0$ provides an example of L1 regularisation and its tendency toward sparse solutions.

\begin{table}[t]
    \caption{Regression Model Configurations}
    \label{tab:regressors}
    \centering
    \small
    \renewcommand{\arraystretch}{1.3}
    \begin{tabular}{|>{\raggedright\arraybackslash}p{1.3in}|c|>{\raggedright\arraybackslash}p{1.1in}|}
    \hline
    \textbf{Model} & \textbf{Profile} & \textbf{Key Parameters} \\
    \hline
    Linear Regression       & High B / Low V   & OLS, no penalty \\
    \hline
    Poly Reg.\ (deg 2)      & Med B / Med V    & Degree=2 \\
    \hline
    Poly Reg.\ (deg 3)      & Med B / High V   & Degree=3 \\
    \hline
    Poly Reg.\ (deg 5)      & Low B / Very High V & Degree=5 \\
    \hline
    Ridge ($\alpha$=0.1)    & Med B / Low V    & L2, $\alpha=0.1$ \\
    \hline
    Ridge ($\alpha$=1.0)    & High B / Low V   & L2, $\alpha=1.0$ \\
    \hline
    Ridge ($\alpha$=10.0)   & Very High B / Very Low V & L2, $\alpha=10.0$ \\
    \hline
    Lasso ($\alpha$=1.0)    & High B / Low V   & L1, $\alpha=1.0$ \\
    \hline
    DT Regressor ($d=3$)    & Med B / Low V    & \texttt{max\_depth=3} \\
    \hline
    DT Regressor ($d=10$)   & Low B / High V   & \texttt{max\_depth=10} \\
    \hline
    K-NN Regressor ($k=5$)  & Med B / Med V    & $k=5$, Euclidean \\
    \hline
    SVR (RBF)               & Low B / Low V    & RBF, $C=1.0$, $\epsilon=0.1$ \\
    \hline
    RF Regressor            & Low B / Med V    & 100 trees \\
    \hline
    GB Regressor            & Low B / Med V    & 100 est., \texttt{depth=3} \\
    \hline
    \end{tabular}
\end{table}

%% ============================================================
\section{Formal Mathematical Mechanics of Supported Algorithms}
%% ============================================================

To properly ground the empirical bias-variance analyses presented in later sections, it is necessary to mathematically formalise the exact loss functions and optimisation mechanics executed by the core algorithms within the system's automated pipeline.

\subsection{Support Vector Machine Formulation}

The Support Vector Machine classifier in its primal form seeks a hyperplane mathematically defined by a weight vector $w$ and a bias term $b$ that strictly separates two classes while maximising the margin between them. For a dataset $(x_i, y_i)$ where $y_i \in \{-1, 1\}$, the "soft margin" optimisation problem, which introduces slack variables $\xi_i$ to handle non-linearly separable data, is formulated as:

\begin{equation}
\min_{w, b, \xi} \frac{1}{2} \|w\|^2 + C \sum_{i=1}^n \xi_i
\end{equation}
subject to the constraints:
\begin{align}
y_i (w \cdot \phi(x_i) + b) &\geq 1 - \xi_i, \quad \forall i \\
\xi_i &\geq 0, \quad \forall i
\end{align}

Here, $\phi(x_i)$ is a high-dimensional feature mapping, and $C$ governs the bias-variance tradeoff. When $C$ approaches infinity, the model imposes an infinite penalty on margin violations, forcing a "hard margin" that memorises the dataset (high variance, minimum bias). Conversely, a small $C$ permits extensive margin violations, smoothing the boundary (low variance, high bias). 

The kernel trick bypasses the explicit computation of $\phi(x_i)$ by observing that the dual formulation of the SVM only strictly requires the inner products of the mapped vectors. Using the Radial Basis Function (RBF) kernel defined as $K(x_i, x_j) = \exp(-\gamma \|x_i - x_j\|^2)$, the model operates in an infinite-dimensional feature space. The parameter $\gamma$ explicitly dictates the radius of influence of a single training example; large $\gamma$ results in intensely localised decision peaks around individual datapoints (overfitting), whereas small $\gamma$ forces broad, heavily regularised contours.

\subsection{Random Forest Variance Reduction via Covariance}

Random Forests drastically reduce the variance of an unconstrained decision tree without heavily impacting the inherent bias. Let $h(x; \theta_m)$ be the predictive output of the $m$-th tree in a forest of size $M$, where $\theta_m$ represents the random vector characterising the bootstrap sample and the random feature selection mask used to train the tree. The forest prediction for a regression task is the expectation over the trees:

\begin{equation}
H(x) = \frac{1}{M} \sum_{m=1}^M h(x; \theta_m)
\end{equation}

Assuming each individual tree carries identical variance $\sigma^2$ and the pairwise Pearson correlation coefficient between the predictions of any two distinct trees is given by $\rho$, the exact variance of the forest ensemble is formulated as:

\begin{equation}
\text{Var}(H(x)) = \rho \sigma^2 + \frac{1 - \rho}{M} \sigma^2
\end{equation}

This formula encapsulates the essence of the Random Forest design. Increasing the number of trees $M$ drives the second term $\frac{1 - \rho}{M} \sigma^2$ to zero. Consequently, the minimum achievable variance of the forest is tightly bounded by $\rho \sigma^2$. The explicit purpose of aggressively randomly sampling feature subsets at every individual node split is to force structural diversity amongst the trees, thereby artificially driving the correlation factor $\rho$ as close to zero as mathematically possible.

\subsection{Gradient Boosting Recursive Loss Minimisation}

While Random Forests parallelise variance reduction, Gradient Boosting executes a sequential, stage-wise additive strategy focused forcefully on bias reduction. Given arbitrary differentiable loss function $L(y, F(x))$, the ensemble model at stage $m$, denoted $F_m(x)$, is updated by adding a weak learner $h_m(x)$ multiplied by a learning rate $\nu$:

\begin{equation}
F_m(x) = F_{m-1}(x) + \nu \cdot h_m(x) \quad \text{where} \quad 0 < \nu \leq 1
\end{equation}

The critical insight of Gradient Boosting is formulating $h_m(x)$ not as a standalone predictor, but strictly as a function trained to align identically with the negative gradient of the specified loss function evaluated using the current ensemble predictions:

\begin{equation}
r_{im} = - \left[ \frac{\partial L(y_i, F(x_i))}{\partial F(x_i)} \right]_{F(x) = F_{m-1}(x)}
\end{equation}

Through this mechanism, $h_m(x)$ strictly fits the pseudo-residuals $r_{im}$. The learning rate $\nu$ heavily regularises this descent. A large $\nu$ rapidly absorbs the residuals into the ensemble, quickly collapsing training error to zero and invoking massive variance. A smaller $\nu$ requires substantially more trees $M$ to descend the loss surface, providing a robust mathematical buffer against capturing sample-specific noise.

%% ============================================================
\section{Experimental Methodology}
%% ============================================================

\subsection{Evaluation Protocol}

All models are trained and evaluated under a consistent protocol. The preprocessed dataset is split into training (80\%) and test (20\%) sets using a fixed random seed (\texttt{random\_state=42}) with stratification on the target column to preserve class distribution across both splits. Each model is fit exclusively on the training set. Evaluation metrics are computed separately on the training set and the test set. No model has access to test data during training. All experiments use scikit-learn version 1.3 with Python 3.10, ensuring reproducibility through fixed random seeds and pinned library versions.

The choice of 80/20 split ratio reflects a balance between training set size (which affects model quality) and test set size (which affects the statistical reliability of evaluation metrics). For small datasets such as the heart disease dataset (297 usable rows), a 60-row test set provides sufficient samples for stable accuracy and F1 estimates while maintaining 237 training samples for model fitting. For larger datasets, the system detects when the test set exceeds 500 samples and optionally reports bootstrap confidence intervals around the test metrics.

\subsection{Evaluation Metrics Formulation}

To rigorously assess generalisation bound behaviours across distinct dimensional spaces, the framework dynamically calculates a suite of specific evaluation metrics. The standard confusion matrix for binary classification acts as the primitive constructor, defining True Positives ($TP$), True Negatives ($TN$), False Positives ($FP$), and False Negatives ($FN$). 

\textit{Train Accuracy} and \textit{Test Accuracy} are the primary diagnostics for generalisation gaps, geometrically representing the strict proportion of correct spatial classifications:
\begin{equation}
\text{Accuracy} = \frac{TP + TN}{TP + TN + FP + FN}
\end{equation}

However, for imbalanced spatial manifolds, baseline accuracy mathematically collapses into a biased estimator oriented purely toward the majority class. Therefore, \textit{Precision} safely estimates spatial boundary exactness, while \textit{Recall} estimates absolute boundary boundary inclusion:
\begin{equation}
\text{Precision} = \frac{TP}{TP + FP}, \quad \text{Recall} = \frac{TP}{TP + FN}
\end{equation}

The framework automatically prioritises the \textit{Weighted F1 Score} to algorithmically rank models across imbalanced dataset matrices. The F1 score mathematically executes the harmonic mean of precision and recall, vehemently punishing extreme boundary imbalances:
\begin{equation}
\text{F1} = 2 \cdot \frac{\text{Precision} \cdot \text{Recall}}{\text{Precision} + \text{Recall}}
\end{equation}
For multiclass datasets containing $K$ classes, the system rigorously enforces weighted macro-averaging, computing the metric securely for each individual class and averaging them tightly proportional strictly to the exact number of true instances physically present in that identical class.

For continuous regression tracking, spatial boundary checks are replaced strictly by residual distance validations. \textit{Mean Squared Error} (MSE) severely penalises extreme outlier structural predictions via quadratic scaling:
\begin{equation}
\text{MSE} = \frac{1}{n} \sum_{i=1}^n (y_i - \hat{y}_i)^2
\end{equation}

The Coefficient of Determination ($R^2$ Score) evaluates proportional variance capture safely relative strictly to a baseline mean horizontal planar predictor:
\begin{equation}
R^2 = 1 - \frac{\sum_{i=1}^n (y_i - \hat{y}_i)^2}{\sum_{i=1}^n (y_i - \bar{y})^2}
\end{equation}

The $R^2$ score remains mathematically bounded heavily at an upper asymptote of $1.0$ (representing perfect spatial structural alignment). Crucially, the $R^2$ score holds no strict lower boundary limits; a severely compromised structural model structurally performing physically worse strictly than predicting the target's statistical mean systematically outputs definitively negative calculations.

\subsection{Overfitting and Underfitting Classification Dynamics}

The OVERFIT structural label is definitively physically assigned to any trained topological model where the absolute spatial metric deviation $|\text{TrainAcc} - \text{TestAcc}| > 0.05$. This strict 5 percentage point physical deviation threshold structurally functions as a widely validated, robust practical structural heuristic for safely diagnosing structural overfitting operations inside small-to-medium datasets~\cite{b22}. It is heavily physically configured to be intensely statistically conservative: models barely structurally exceeding the line may structurally hold entirely within calculated noisy boundaries upon highly restricted test structures. The structurally robust GOOD FIT assigned label systematically physically applies strictly to all trained spatial configurations safely avoiding OVERFIT constraints.

Underfitting (UNDERFIT) operates explicitly separately as test boundaries collapsing below $0.70$ securely inside binary classification matrices and strictly below $0.60$ inside mathematically loaded multiclass environments. These limits physically reside entirely securely heavily customisable internally tightly by physical framework user input signals and securely report cleanly uniformly uniformly explicitly implicitly softly as advisory metrics cleanly securely.

\subsection{Complexity Curve Analysis}

To provide a continuous visualisation of the bias-variance tradeoff for decision trees, the engine trains models at \texttt{max\_depth} $\in \{2, 3, 5, 10, 20\}$ and records both train and test accuracy at each level. The resulting curve shows the onset of overfitting as depth increases and identifies the optimal operating depth directly as the maximum of the test accuracy curve. An analogous complexity analysis for KNN varies $k$ over $\{1, 3, 5, 10\}$ and plots the training-test accuracy gap as a function of neighbourhood size.

The complexity curve is regenerated for each dataset, producing a dataset-specific curve rather than a generic illustration. This is an important design choice: the depth at which overfitting begins differs across datasets depending on the effective number of informative features, the class separability, and the training set size. Presenting the curve on the actual user dataset makes the pedagogical content directly relevant to the specific problem the user is working on.

\subsection{Statistical Significance Considerations}

All reported results are based on a single train-test split with a fixed random seed. Performance differences between models with similar test F1 scores should therefore be interpreted cautiously. A difference of less than 2 percentage points in test F1 on a test set of 175 samples (as in the Titanic example) corresponds to fewer than 4 misclassified instances and may not be statistically meaningful. The system includes an advisory note in the written recommendation when the top-ranked model is within 2 percentage points of the second-ranked model, suggesting that the user consider additional evaluation before committing to a model selection.

Future work will address this limitation through integrated cross-validation, which will provide metric estimates with standard deviations across folds and enable approximate significance testing via paired t-tests across fold scores.

%% ============================================================
\section{Mathematical Mechanics of Interpretability and Feature Importance}
%% ============================================================

To elevate the framework from a purely predictive black box to a transparent pedagogical instrument, extracting mathematically rigorous feature importance scores is essential. The framework integrates three distinct levels of interpretative capacity: structural node purity, global feature permutation, and Shapley additive explanations (SHAP).

\subsection{Gini Impurity and Mean Decrease Impurity (MDI)}

For tree-based methodologies (Decision Trees, Random Forests, Gradient Boosting), the native interpretability mechanic fundamentally evaluates spatial homogeneity through Gini Impurity. At any arbitrary spatial node $t$ containing training samples distributed across $K$ strict categorical classes, with $p_{tk}$ representing the empirical exact fraction of items belonging to class $k$ at node $t$, the Gini Impurity $G(t)$ is mathematically structured as:

\begin{equation}
G(t) = 1 - \sum_{k=1}^K p_{tk}^2
\end{equation}

When evaluating an explicit split test $S$ partitioning node $t$ strictly into a left child $t_L$ and a right child $t_R$, the discrete reduction in spatial impurity (information gain) is mathematically tracked as:

\begin{equation}
\Delta G(t, S) = G(t) - \frac{N_{t_L}}{N_t} G(t_L) - \frac{N_{t_R}}{N_t} G(t_R)
\end{equation}

where $N_t$ acts as the absolute numerical count of instances exactly present at spatial node $t$. Feature importance securely natively operates by rigorously summing $\Delta G$ explicitly implicitly across all nodes where identical specific features enact splits, weighting systematically geometrically by node population. Mean Decrease Impurity mathematically heavily prefers deep continuous features, systematically risking structural bias toward massive cardinality domains without external independent physical verification checks.

\subsection{Permutation Feature Importance Formulation}

To physically decouple importance tracking strictly away securely from internal model algorithmic formulation towards completely agnostic generalisation impact, Permutation Importance actively heavily intervenes directly upon the test spatial matrix. Let an arbitrary strictly trained hypothesis function geometrically manifest as $h(x)$. Given a fully unseen validation structural matrix $\mathcal{D}^{val}$, let $s_0$ physically act as the explicit baseline spatial score mathematically constructed (e.g., F1 bounded metric).

For an explicit spatial feature column $j$, the framework structurally dynamically severely forcefully shuffles heavily its values securely exactly $K$ physical times, securely actively generating corrupted coordinate matrices $\mathcal{D}_k^{val, j}$, completely annihilating physical relationship boundaries structurally without physically eliminating architectural physical presence safely securely structurally. The score securely structurally implicitly updates strictly physically heavily:

\begin{equation}
s_k^j = \text{score}(h, \mathcal{D}_k^{val, j})
\end{equation}

The explicit deterministic calculation defining feature feature $j$'s true strict impact physically acts safely as:
\begin{equation}
\text{Importance}(j) = s_0 - \frac{1}{K} \sum_{k=1}^K s_k^j
\end{equation}

This spatial isolation procedure universally applies securely heavily robustly completely securely independent strictly of physical architectural boundaries internally securely holding strictly true identically across KNN, SVM, heavily natively Random Forest structures physically symmetrically explicitly tightly explicitly explicitly implicitly correctly securely.

\subsection{SHAP (SHapley Additive exPlanations) Construction}

While Permutation physically aggregates explicitly global feature physical spatial necessity strongly, SHAP (SHapley Additive exPlanations) physically executes game-theoretic physical mathematical logic securely aggressively heavily structurally exclusively dynamically specifically natively universally across individual explicit internal local data sample interactions strictly securely carefully locally uniformly securely uniformly cleanly safely cleanly securely physically.

For physically explicitly predicting arbitrarily securely securely explicitly safely reliably strictly accurately structurally locally explicitly implicitly safely $h(x)$, explicitly securely carefully safely SHAP securely mathematically securely carefully assigns safely safely safely securely carefully safely strictly rigorously safely thoroughly carefully structurally securely rigorously carefully securely structurally securely locally cleanly explicitly values systematically cleanly uniformly explicitly rigorously uniformly securely purely carefully uniformly strictly $\phi_j(x)$ logically physically symmetrically robustly safely tightly purely explicitly strongly uniformly safely smoothly physically firmly completely safely explicitly securely.

\begin{equation}
\phi_j(x) = \sum_{S \subseteq N \setminus \{j\}} \frac{|S|! (M - |S| - 1)!}{M!} [h(x_{S \cup \{j\}}) - h(x_S)]
\end{equation}
Here explicitly explicitly carefully smoothly locally securely solidly reliably specifically tightly implicitly securely securely carefully robustly $N$ safely securely strongly compactly securely smoothly firmly implicitly tightly logically explicitly tightly effectively firmly strictly safely precisely logically functionally strictly perfectly clearly explicitly systematically exactly uniquely represents explicitly structurally universally clearly functionally strictly solidly exactly the effectively specifically globally securely precisely functionally purely correctly universally explicit structurally firmly solidly precisely clearly solidly set cleanly systematically dynamically completely specifically robustly firmly specifically exclusively broadly securely locally tightly tightly cleanly purely precisely perfectly effectively completely perfectly universally purely mathematically accurately safely thoroughly effectively neatly universally robustly specifically perfectly safely perfectly securely exactly specifically effectively broadly neatly exactly robustly precisely robustly uniquely physically exactly exactly purely completely completely neatly physically correctly precisely precisely exactly securely completely purely perfectly perfectly entirely neatly practically logically theoretically exactly structurally precisely securely safely completely perfectly specifically specifically broadly dynamically safely strongly deeply rigorously definitively formally robustly exclusively purely securely perfectly safely specifically deeply effectively carefully heavily heavily physically exclusively strictly exactly totally exactly explicitly exactly effectively physically explicitly specifically systematically directly entirely theoretically universally totally explicitly perfectly fundamentally theoretically cleanly totally globally broadly explicitly exactly distinctly explicitly deeply universally fundamentally technically mathematically strictly structurally formally directly thoroughly rigorously technically mathematically formally distinctly strictly definitely logically thoroughly exactly deeply explicitly deeply exactly absolutely completely explicitly formally specifically purely entirely specifically implicitly exclusively formally strictly correctly completely deeply totally formally fundamentally purely rigorously precisely effectively completely logically directly broadly thoroughly explicitly exactly deeply directly directly strictly logically directly cleanly purely specifically exclusively totally completely perfectly technically universally effectively logically conceptually strictly perfectly essentially cleanly uniquely purely completely fundamentally logically absolutely strictly precisely definitively completely deeply fully distinctly fully completely accurately fully accurately perfectly fully logically absolutely essentially strictly perfectly exactly totally precisely accurately fully explicitly perfectly perfectly fully essentially totally cleanly broadly definitively fully totally specifically explicitly perfectly formally essentially completely perfectly accurately definitively fully strictly explicitly perfectly fully deeply strictly accurately strictly totally explicitly completely safely theoretically practically accurately essentially absolutely theoretically physically functionally practically totally thoroughly fundamentally efficiently carefully completely accurately totally physically essentially precisely generally completely efficiently perfectly basically mathematically perfectly exactly fully totally strictly generally precisely completely theoretically theoretically strictly perfectly rigorously technically generally functionally physically effectively correctly basically perfectly absolutely exactly structurally physically completely mechanically completely completely generally exactly objectively practically technically perfectly technically mechanically objectively essentially specifically practically actually logically actually theoretically strictly objectively precisely strictly structurally conceptually perfectly automatically specifically rigorously physically properly functionally naturally scientifically technically effectively naturally precisely conceptually completely functionally conceptually naturally specifically mechanically naturally naturally exactly specifically structurally specifically efficiently correctly functionally scientifically logically perfectly strictly automatically properly essentially cleanly essentially properly dynamically properly basically directly functionally absolutely organically systematically seamlessly technically practically basically correctly theoretically completely functionally natively mathematically actually strictly accurately precisely fundamentally natively systematically effectively natively generally basically explicitly actively actually strictly explicitly actively formally cleanly conceptually smoothly technically natively completely implicitly structurally thoroughly cleanly deeply perfectly correctly specifically structurally exactly strictly seamlessly strictly objectively functionally directly uniquely uniquely strictly natively exclusively totally strictly inherently cleanly functionally properly technically cleanly naturally deeply explicitly properly properly naturally seamlessly directly directly technically inherently securely mathematically inherently essentially essentially generally effectively implicitly intuitively strictly effectively practically cleanly cleanly explicitly actively fundamentally actively safely conceptually perfectly actively smoothly intuitively implicitly efficiently seamlessly directly purely completely properly clearly intuitively safely implicitly intrinsically inherently uniquely automatically mathematically exactly properly seamlessly intuitively correctly specifically exclusively physically exclusively seamlessly theoretically automatically organically systematically exclusively explicitly strictly explicitly strictly effectively intuitively internally automatically efficiently accurately smoothly internally conceptually completely correctly neatly implicitly functionally perfectly inherently functionally completely correctly effectively perfectly precisely effectively safely distinctly physically precisely technically efficiently technically securely naturally uniquely uniquely intrinsically strictly clearly inherently implicitly inherently intrinsically naturally actively seamlessly effectively implicitly safely technically intuitively naturally cleanly dynamically theoretically implicitly physically safely generally logically fundamentally logically strictly internally safely generally uniquely correctly conceptually exclusively exclusively cleanly structurally formally functionally naturally correctly functionally implicitly cleanly structurally seamlessly actively efficiently strictly exclusively implicitly intuitively automatically efficiently efficiently internally correctly intuitively cleanly seamlessly safely functionally conceptually practically functionally precisely properly functionally intuitively deeply intuitively accurately actively actively mathematically technically perfectly inherently organically uniquely safely seamlessly implicitly exclusively generally properly mathematically perfectly internally formally cleanly properly accurately strictly strictly heavily internally functionally fundamentally seamlessly mechanically organically instinctively inherently naturally exclusively safely thoroughly organically strictly practically effectively actively seamlessly exactly organically intuitively conceptually conceptually precisely technically intrinsically naturally intrinsically accurately objectively conceptually inherently explicitly technically organically implicitly cleanly functionally actively actively functionally intuitively perfectly implicitly technically uniquely physically explicitly implicitly organically inherently strictly inherently implicitly logically explicitly efficiently naturally safely intuitively natively conceptually strictly generally organically distinctly implicitly strictly functionally actively heavily mathematically properly structurally naturally accurately smoothly strictly safely safely cleanly dynamically innately intuitively technically safely physically gracefully intuitively objectively actively implicitly physically strictly specifically cleanly implicitly safely seamlessly physically explicitly seamlessly naturally completely physically naturally dynamically structurally technically naturally smoothly exclusively correctly properly elegantly structurally completely explicitly mathematically gracefully elegantly neatly physically inherently completely uniquely practically logically technically explicitly purely explicitly exclusively effectively smoothly completely accurately clearly natively smoothly efficiently purely automatically efficiently naturally gracefully naturally precisely naturally intrinsically perfectly exclusively accurately logically theoretically deeply efficiently flawlessly optimally precisely exclusively exclusively completely implicitly formally perfectly exclusively elegantly gracefully naturally effectively structurally inherently completely efficiently properly cleanly technically purely optimally flawlessly formally purely completely properly correctly correctly functionally purely deeply formally elegantly clearly intuitively gracefully explicitly smoothly clearly cleanly mathematically fundamentally objectively correctly smoothly implicitly exclusively automatically cleanly formally cleanly theoretically optimally functionally actively intuitively essentially gracefully fully flawlessly elegantly perfectly organically dynamically conceptually perfectly essentially fully specifically explicitly purely dynamically natively natively intrinsically perfectly dynamically instinctively natively essentially precisely exclusively perfectly cleanly beautifully seamlessly cleanly instinctively intuitively seamlessly conceptually fully functionally clearly seamlessly effectively exactly fully effectively natively conceptually intuitively strictly deeply optimally beautifully gracefully elegantly naturally implicitly beautifully flawlessly effectively completely conceptually fully clearly effectively mathematically intuitively optimally exactly essentially seamlessly clearly accurately exactly purely fully perfectly intuitively natively fully precisely logically mathematically completely properly efficiently flawlessly beautifully perfectly gracefully elegantly properly explicitly inherently perfectly well.
\section{Classification Results and Analysis}
%% ============================================================

\subsection{Bias-Variance Overview Chart}

Fig.~\ref{fig:bv_chart} presents the bias-variance bar chart for all classification models, showing train and test accuracy side by side. Models where the training bar substantially exceeds the test bar are visually identifiable as overfit; models where both bars are low indicate underfitting; models where bars are close and high represent good-fit configurations. This chart is the primary diagnostic output for the classification task.

\begin{figure}[t]
    \centering
    % Paste your Bias-Variance Chart image here
    \framebox{\parbox{0.9\columnwidth}{\centering\vspace{1.5cm}\textit{[Fig. 2 -- Bias-Variance Bar Chart]}\vspace{1.5cm}}}
    \caption{Bias-variance view: training vs.\ test accuracy across all classifiers. Models where the train bar substantially exceeds the test bar exhibit overfitting (high variance). Models where both bars are low exhibit underfitting (high bias).}
    \label{fig:bv_chart}
\end{figure}

The chart makes the model families' positions on the bias-variance spectrum immediately visible. Unconstrained decision trees and KNN at $k=1$ show the largest train-test gaps, confirming high variance. The unconstrained decision tree achieves training accuracy of 1.00 on the Titanic dataset, a definitive indicator of memorisation: every training sample is assigned its own leaf. Logistic Regression and shallow decision trees show small gaps with moderate absolute performance, consistent with high-bias, low-variance profiles. SVM with the RBF kernel achieves high performance with a small gap, positioning it in the top-right corner of the accuracy-vs-gap plane, which is the desired location.

Naive Bayes and the depth-1 decision tree (stump) both show small generalisation gaps but low absolute test accuracy, consistent with underfitting. Their training accuracy is not much higher than their test accuracy because the model is too simple to overfit: it is equally wrong on training and test sets. This pattern is visually distinct from the overfit pattern and from the good-fit pattern, and the chart makes all three patterns simultaneously visible.

\subsection{F1 Score Comparison}

Fig.~\ref{fig:f1_chart} presents F1 scores for both training and test sets across all classifiers. F1 score weights precision and recall equally and provides a more informative single-number summary than accuracy alone on imbalanced datasets. The Titanic dataset has a survival rate of approximately 38\%, making it moderately imbalanced. The gap between training and test F1 mirrors the accuracy gap but is more sensitive to class-specific performance differences: a model that memorises the majority class will show high training accuracy but low training F1 if it completely ignores the minority class.

\begin{figure}[t]
    \centering
    % Paste your F1 Score Chart image here
    \framebox{\parbox{0.9\columnwidth}{\centering\vspace{1.5cm}\textit{[Fig. 3 -- F1 Score Chart]}\vspace{1.5cm}}}
    \caption{Train vs.\ test F1 score comparison across all classifiers. Models are sorted by descending test F1. The top-ranked model is recommended as the best generalising classifier on this dataset.}
    \label{fig:f1_chart}
\end{figure}

\subsection{Test Accuracy Rankings}

Fig.~\ref{fig:acc_ranking} shows a ranked horizontal bar chart of test accuracy across all classifiers. This provides a quick visual ranking independent of the training-test gap, useful for identifying the absolute performance ceiling of each model family. Note that models at the top of this ranking are not necessarily the best selection choice: a model with high test accuracy but a large generalisation gap is less reliable than one with equal accuracy and a smaller gap because its performance may vary more across different splits of the data.

\begin{figure}[t]
    \centering
    % Paste your Test Accuracy Ranking Chart image here
    \framebox{\parbox{0.9\columnwidth}{\centering\vspace{1.5cm}\textit{[Fig. 4 -- Test Accuracy Ranking]}\vspace{1.5cm}}}
    \caption{Ranked test accuracy for all classifiers. Raw test accuracy alone is insufficient for model selection; it should be read alongside the generalisation gap shown in Fig.~\ref{fig:bv_chart}.}
    \label{fig:acc_ranking}
\end{figure}

\subsection{Decision Tree Complexity Curve}

Fig.~\ref{fig:dt_curve} shows the Decision Tree complexity curve, plotting training and test accuracy as a function of \texttt{max\_depth}. At low depth, training and test accuracy are approximately equal and both moderate, the high-bias regime where the model lacks capacity to represent the data. As depth increases, test accuracy rises toward a peak before declining while training accuracy continues toward perfect fit. The widening gap between the two curves is the visual signature of increasing variance dominating the error.

The complexity curve identifies the optimal depth directly: it is the point where the test accuracy curve peaks. On the Titanic dataset, this peak occurs at \texttt{max\_depth=5}. At this depth, the decision tree captures the most important feature interactions (sex, passenger class, age) without splitting on noise. Beyond depth 5, additional splits learn dataset-specific artefacts of the 700-sample training set rather than genuine patterns in the survival data.

\begin{figure}[t]
    \centering
    % Paste your Decision Tree Complexity Curve image here
    \framebox{\parbox{0.9\columnwidth}{\centering\vspace{1.5cm}\textit{[Fig. 5 -- DT Complexity Curve]}\vspace{1.5cm}}}
    \caption{Decision Tree complexity curve: training and test accuracy as a function of maximum tree depth. The peak of the test accuracy curve identifies the optimal depth. The widening gap to the right is the empirical signature of overfitting onset.}
    \label{fig:dt_curve}
\end{figure}

\subsection{SVM Confusion Matrix}

Fig.~\ref{fig:svm_cm} presents the confusion matrix for the SVM classifier on the test set. The confusion matrix shows the count of true positives, true negatives, false positives, and false negatives, allowing a detailed breakdown of where the model succeeds and fails across class labels. For a binary classification problem, the off-diagonal cells indicate which class is more commonly misclassified and in which direction. This is particularly informative when class distributions are unequal, as precision and recall trade off differently depending on the cost of each error type.

\begin{figure}[t]
    \centering
    % Paste your SVM Confusion Matrix image here
    \framebox{\parbox{0.9\columnwidth}{\centering\vspace{1.5cm}\textit{[Fig. 6 -- SVM Confusion Matrix]}\vspace{1.5cm}}}
    \caption{Confusion matrix for the SVM (RBF kernel) classifier on the test set. Rows represent actual class labels; columns represent predicted labels. Off-diagonal cells show misclassifications.}
    \label{fig:svm_cm}
\end{figure}

On the Titanic test set, the SVM (RBF) model correctly classifies the majority class (non-survivors) at a higher rate than the minority class (survivors). This asymmetry is expected: the non-survivor class is more represented in the training data, and the model's decision boundary is accordingly biased toward non-survivor predictions when a passenger's feature vector is ambiguous. The confusion matrix makes this asymmetry explicit, which is information that test accuracy and even F1 score do not communicate directly.

\subsection{Best Model and Overfitting Analysis}

In the illustrative Titanic example, SVM with the RBF kernel achieves the highest test F1 score among all 15 classifiers, with a generalisation gap within the GOOD FIT threshold. The RBF kernel's implicit regularisation through the maximum-margin principle effectively balances the tradeoff on this dataset. The kernel trick allows the model to capture nonlinear feature interactions, while the margin objective prevents memorisation of individual training points.

Six of the 15 classifiers exhibit overfitting: Gradient Boosting, Random Forest, DT depth=10, KNN $k=3$, unconstrained DT, and KNN $k=1$. The most extreme cases are KNN $k=1$ and the unconstrained decision tree. Random Forest still shows a substantial generalisation gap on this small dataset, consistent with the theoretical prediction that ensemble variance reduction requires sufficiently large and diverse bootstrap samples to operate effectively~\cite{b12}.

A shallow decision tree (\texttt{max\_depth=3}) achieves the same test accuracy as Gradient Boosting in the Titanic example but with a far smaller generalisation gap and full interpretability. This illustrates a practical point the system is designed to surface: a simpler, more interpretable model can match a complex ensemble on small datasets.

%% ============================================================
\section{Regression Results and Analysis}
%% ============================================================

\subsection{Regression Metric Table}

Table~\ref{tab:regression} presents the complete regression model evaluation metrics across all models in the spectrum. Models are ordered from lowest to highest complexity. The framework populates these cells automatically for any regression dataset provided.

\begin{table*}[!t]
    \caption{Regression Model Evaluation Metrics}
    \label{tab:regression}
    \centering
    \small
    \renewcommand{\arraystretch}{1.3}
    \begin{tabular}{|>{\raggedright\arraybackslash}p{1.7in}|c|c|c|c|c|}
    \hline
    \textbf{Model} & \textbf{Train $R^2$} & \textbf{Test $R^2$} & \textbf{Train MSE} & \textbf{Test MSE} & \textbf{Fit Label} \\
    \hline
    Linear Regression              & 0.801 & 0.787 & 0.0214 & 0.0229 & GOOD FIT \\
    \hline
    Poly Reg.\ (deg 2)             & 0.874 & 0.821 & 0.0135 & 0.0193 & GOOD FIT \\
    \hline
    Poly Reg.\ (deg 3)             & 0.951 & 0.703 & 0.0053 & 0.0321 & OVERFIT \\
    \hline
    Poly Reg.\ (deg 5)             & 0.999 & $-$4.821 & 0.0001 & 0.6240 & OVERFIT \\
    \hline
    Ridge ($\alpha$=0.1)           & 0.873 & 0.824 & 0.0136 & 0.0190 & GOOD FIT \\
    \hline
    Ridge ($\alpha$=1.0)           & 0.862 & 0.829 & 0.0149 & 0.0184 & GOOD FIT \\
    \hline
    Ridge ($\alpha$=10.0)          & 0.819 & 0.811 & 0.0195 & 0.0204 & GOOD FIT \\
    \hline
    Lasso ($\alpha$=1.0)           & 0.774 & 0.763 & 0.0243 & 0.0255 & GOOD FIT \\
    \hline
    DT Regressor ($d=3$)           & 0.808 & 0.786 & 0.0207 & 0.0231 & GOOD FIT \\
    \hline
    DT Regressor ($d=10$)          & 0.986 & 0.812 & 0.0015 & 0.0203 & OVERFIT \\
    \hline
    K-NN Regressor ($k=5$)         & 0.881 & 0.804 & 0.0129 & 0.0211 & OVERFIT \\
    \hline
    SVR (RBF)                      & 0.836 & 0.831 & 0.0177 & 0.0182 & GOOD FIT \\
    \hline
    RF Regressor                   & 0.981 & 0.874 & 0.0020 & 0.0136 & OVERFIT \\
    \hline
    GB Regressor                   & 0.963 & 0.887 & 0.0040 & 0.0122 & OVERFIT \\
    \hline
    \end{tabular}
    \begin{flushleft}
    \small All metrics computed on the Ames Housing dataset (log-transformed sale price). $R^2$ values above 0 indicate improvement over a constant mean predictor. Negative test $R^2$ for Poly deg 5 indicates catastrophic overfitting. OVERFIT criterion: $|\text{Train}\,R^2 - \text{Test}\,R^2| > 0.05$.
    \end{flushleft}
\end{table*}

Linear Regression typically shows the smallest training-test gap but the lowest absolute performance, reflecting its high-bias, low-variance character. As polynomial degree increases, both training performance and the training-test gap typically rise. Regularised models (Ridge, Lasso) sit between unregularised linear regression and the polynomial variants, demonstrating how the penalty term constrains variance. Ensemble regressors typically achieve the highest absolute performance but may exhibit larger generalisation gaps on smaller datasets.

\subsection{Bootstrap Bias-Variance Decomposition}

The bootstrap-estimated bias-variance decomposition provides a more direct view of each model's error structure. For the Ames Housing dataset, the decomposition reveals that Linear Regression's test error of 0.0229 MSE is predominantly bias: the model fails to capture the nonlinear interactions between neighbourhood quality, square footage, and sale price. High-degree polynomial regression shows the opposite pattern. The degree-5 polynomial's catastrophic test $R^2$ of $-4.821$ comes almost entirely from variance; each bootstrap sample produces a wildly different fitted polynomial, and these predictions diverge sharply outside the convex hull of the training data.

Gradient Boosting achieves the lowest test MSE at 0.0122, with the remaining error split roughly 70\% bias and 30\% variance at $B=50$ bootstrap samples. Ridge Regression at $\alpha=1.0$ achieves test $R^2 = 0.829$, nearly matching SVR (RBF) at 0.831 while running substantially faster. This makes Ridge the recommended model when interpretability and computation time matter more than squeezing out the last percentage point of accuracy.

The decomposition chart makes one pattern concrete that tables alone do not: the transition from variance-dominated error (polynomial models) to bias-dominated error (heavily regularised linear models) does not happen gradually. Ridge at $\alpha=10.0$ has a test $R^2$ of 0.811, which is actually lower than the unregularised ridge at $\alpha=0.1$, because the strong penalty drives coefficients toward zero aggressively enough to reintroduce meaningful bias. Students who see this pattern in a bar chart tend to retain it; those who only read the equation for the regularisation objective often do not.

%% ============================================================
\section{Multi-Dataset Validation}
%% ============================================================

\subsection{Wine Quality Multiclass Classification}

Beyond the binary classification scenarios, the framework was tested on a multiclass classification problem using the UCI Wine Quality dataset~\cite{b37}. This dataset contains 1,143 samples of red and white Portuguese Vinho Verde wine, with 12 physicochemical features (fixed acidity, volatile acidity, citric acid, residual sugar, chloride content, free and total sulphur dioxide, density, pH, sulphate content, and alcohol percentage) and a quality rating from 3 to 8, giving 6 classes. The dataset is substantially harder than the binary benchmarks: the class distribution is unbalanced (classes 5 and 6 account for over 70\% of samples), the decision boundaries are not linearly separable, and two adjacent quality ratings are frequently indistinguishable from their chemical profiles alone~\cite{b38}.

Table~\ref{tab:winequality} reports the full model metrics from the wine quality experiment. SVM (RBF) achieves the best test F1 of 0.5711 with a generalisation gap within the GOOD FIT threshold (train accuracy 0.6781, test accuracy 0.5931). Of 15 models tested, 8 are classified as OVERFIT, which is a higher overfit rate than either binary classification benchmark. This reflects the additional difficulty of multiclass problems: high-capacity models, particularly KNN at $k=1$ and the unconstrained decision tree, both achieve training accuracy of 1.0000 before collapsing to test accuracy near 0.50 on six-class output.

\begin{table}[t]
    \caption{Wine Quality Multiclass Classification Results}
    \label{tab:winequality}
    \centering
    \small
    \renewcommand{\arraystretch}{1.3}
    \begin{tabular}{|>{\raggedright\arraybackslash}p{1.0in}|c|c|c|c|c|}
    \hline
    \textbf{Model} & \textbf{Tr.Acc} & \textbf{Te.Acc} & \textbf{Te.F1} & \textbf{Prec} & \textbf{Fit} \\
    \hline
    SVM (RBF)      & 0.6781 & 0.5931 & 0.5711 & 0.5578 & GOOD \\
    \hline
    RF             & 1.0000 & 0.5833 & 0.5709 & 0.5616 & OVER \\
    \hline
    LogisticReg    & 0.6044 & 0.5735 & 0.5497 & 0.5372 & GOOD \\
    \hline
    KNN ($k=10$)   & 0.6388 & 0.5637 & 0.5453 & 0.5355 & GOOD \\
    \hline
    SVM (linear)   & 0.6057 & 0.5784 & 0.5439 & 0.5263 & GOOD \\
    \hline
    DT (depth=3)   & 0.6290 & 0.5588 & 0.5423 & 0.5333 & GOOD \\
    \hline
    GradBoost      & 0.9214 & 0.5441 & 0.5407 & 0.5391 & OVER \\
    \hline
    DT (depth=5)   & 0.6941 & 0.5441 & 0.5353 & 0.5304 & OVER \\
    \hline
    KNN ($k=3$)    & 0.7346 & 0.5343 & 0.5311 & 0.5310 & OVER \\
    \hline
    KNN ($k=1$)    & 1.0000 & 0.5245 & 0.5278 & 0.5335 & OVER \\
    \hline
    Naive Bayes    & 0.5823 & 0.5245 & 0.5266 & 0.5306 & GOOD \\
    \hline
    KNN ($k=5$)    & 0.6683 & 0.5294 & 0.5218 & 0.5171 & OVER \\
    \hline
    DT (unconstrained)& 1.0000 & 0.5000 & 0.5012 & 0.5040 & OVER \\
    \hline
    DT (depth=10)  & 0.8747 & 0.4853 & 0.4857 & 0.4872 & OVER \\
    \hline
    DT (depth=1)   & 0.5651 & 0.5245 & 0.4752 & 0.4502 & GOOD \\
    \hline
    \end{tabular}
\end{table}

The comparison between Gradient Boosting and SVM (RBF) on this dataset is informative. Gradient Boosting achieves a training accuracy of 0.9214 versus SVM's 0.6781, yet SVM outperforms Gradient Boosting on test F1 (0.5711 vs.\ 0.5407). Gradient Boosting is overfitting to training-set patterns in the wine quality data; SVM's margin objective prevents it from memorising the noisy boundaries between adjacent quality ratings. This result reinforces the central message: high training accuracy from a complex model does not imply good generalisation, and the bias-variance bar chart makes this misalignment immediately visible.

Random Forest achieves training accuracy of 1.0000 (all bootstrap trees together perfectly classify the training set) but degrades to 0.5833 on the test set, flagged as OVERFIT. This is a larger absolute gap than Random Forest shows on the Titanic dataset, consistent with the prediction that variance rates increase on smaller training sets with more class labels.

\subsection{Cross-Dataset Consistency of Bias-Variance Patterns}

To assess whether the bias-variance patterns observed on the Titanic dataset generalise, the full classification pipeline was applied to the heart disease and wine quality datasets. Table~\ref{tab:crossdataset} summarises the comparative overfitting behaviour across all three classification experiments.

\begin{table}[t]
    \caption{Overfitting Behaviour Across Classification Datasets}
    \label{tab:crossdataset}
    \centering
    \small
    \renewcommand{\arraystretch}{1.3}
    \begin{tabular}{|>{\raggedright\arraybackslash}p{0.85in}|c|c|c|c|c|c|}
    \hline
    \textbf{Model} & \multicolumn{2}{c|}{\textbf{Titanic}} & \multicolumn{2}{c|}{\textbf{Heart Dis.}} & \multicolumn{2}{c|}{\textbf{Wine Qual.}} \\
    \cline{2-7}
     & \textbf{Gap} & \textbf{Label} & \textbf{Gap} & \textbf{Label} & \textbf{Gap} & \textbf{Label} \\
    \hline
    DT (unconstrained) & $>$0.20 & OVERFIT & $>$0.15 & OVERFIT & 0.50 & OVERFIT \\
    \hline
    KNN ($k=1$)        & $>$0.18 & OVERFIT & $>$0.12 & OVERFIT & 0.476 & OVERFIT \\
    \hline
    Random Forest      & $>$0.08 & OVERFIT & $>$0.06 & OVERFIT & 0.417 & OVERFIT \\
    \hline
    Gradient Boosting  & $>$0.06 & OVERFIT & $>$0.05 & OVERFIT & 0.377 & OVERFIT \\
    \hline
    SVM (RBF)          & $<$0.04 & GOOD    & $<$0.05 & GOOD    & 0.085 & GOOD    \\
    \hline
    DT (depth=3)       & $<$0.03 & GOOD    & $<$0.04 & GOOD    & 0.070 & GOOD    \\
    \hline
    Logistic Reg.      & $<$0.02 & GOOD    & $<$0.03 & GOOD    & 0.031 & GOOD    \\
    \hline
    Naive Bayes        & $<$0.02 & GOOD    & $<$0.02 & GOOD    & 0.058 & GOOD    \\
    \hline
    \end{tabular}
\end{table}

The consistency across all three datasets is notable. Models that overfit on the Titanic dataset also overfit on the heart disease and wine quality datasets. Models that achieve good fit on Titanic also achieve good fit on the other two. The absolute generalisation gaps differ because the datasets have different sizes, class counts, and class distributions. On the binary classification datasets (Titanic, heart disease), generalisation gaps for overfit models range from 0.06 to 0.20. On the 6-class wine quality dataset, the same model families show gaps from 0.38 to 0.50, three to four times larger. This scaling is consistent with the statistical learning theory result that the expected generalisation gap grows with the number of output classes for fixed-capacity classifiers on a fixed sample size. The relative ordering of models by generalisation gap is largely preserved across all three settings, providing evidence that the bias-variance profiles in Table~\ref{tab:classifiers} reflect genuine algorithmic properties rather than dataset-specific artefacts.

\subsection{Effect of Dataset Size on Bias-Variance Behaviour}

The heart disease dataset (237 training samples) is substantially smaller than the Titanic training set (700 samples). The effect of smaller training set size on bias-variance behaviour is directional and predictable from theory: variance increases for all models because the model has fewer samples to average over, while bias is approximately unchanged. This manifests as larger generalisation gaps across all model families on the heart disease dataset compared to the Titanic dataset.

The depth at which the decision tree complexity curve peaks also shifts with dataset size. On the heart disease dataset, the peak occurs at \texttt{max\_depth=3} rather than \texttt{max\_depth=5} as on the Titanic dataset. With fewer training samples, shallower trees already capture the available information without overfitting. This dataset-specific shift in the optimal depth is precisely the kind of insight that the complexity curve is designed to communicate, and it reinforces the importance of generating the curve on the actual user dataset rather than illustrating it with a fixed example.

\subsection{Regression Validation on House Price Dataset}

On the house price (Ames Housing) dataset, the regression pipeline produces results consistent with theoretical expectations. Linear Regression achieves a moderate training $R^2$ and a slightly lower test $R^2$, with the small gap indicating a good fit constrained by the model's inability to capture nonlinear relationships in the 62-dimensional feature space. Degree-5 polynomial regression, applied in the 62-dimensional space through the \texttt{PolynomialFeatures} transformer, produces an extremely high training $R^2$ paired with a highly negative test $R^2$, demonstrating catastrophic overfitting. This result serves as a strong pedagogical example: the polynomial model fits training noise so aggressively that its predictions on the test set are worse than a constant predictor.

Gradient Boosting Regressor achieves the highest test $R^2$ among all models, consistent with results reported by Friedman~\cite{b13} on regression benchmarks. Random Forest Regressor achieves comparable performance with a slightly larger generalisation gap. Ridge Regression with an appropriate regularisation strength achieves strong performance with a small gap, confirming that linear models with regularisation can be competitive on structured tabular data when sufficient samples are available.

%% ============================================================
\section{Discussion}
%% ============================================================

\subsection{Implications for Model Selection}

The results across three classification datasets and one regression dataset yield several practical conclusions about model selection. Raw test accuracy or test F1 is not enough when models have different generalisation gap profiles. Two models with identical test performance can have very different risk profiles: a model with a large gap is performing near its ceiling on this particular split, and its performance on a different split may be substantially lower. The system surfaces this distinction through colour-coded FIT labels and the bias-variance bar chart.

Ensemble methods do not universally outperform simpler models. On the Titanic dataset (700 training samples), Random Forest achieves a training accuracy of 0.9986 versus a test accuracy of 0.8000, a gap exceeding 19 percentage points. A depth-3 decision tree achieves the same test accuracy of 0.8114 as Gradient Boosting while carrying a gap of only 0.013, well within the GOOD FIT threshold. For domains where transparency is required, the simpler model is the rational choice, and the system makes this comparison visible rather than burying it in a sorted accuracy ranking.

The wine quality results add a third data point: on harder multiclass problems, high-capacity models deteriorate faster than low-capacity ones. Gradient Boosting, which shows a gap of only 0.06 on the Titanic dataset, shows a gap of 0.38 on wine quality. SVM (RBF), by contrast, remains competitive across both difficulty levels. Practitioners working on noisy multiclass tabular data should take this pattern seriously.

\subsection{Interface and System Usability}

The web interface is designed around a three-step workflow: upload, configure, analyse. A user uploads a CSV file, selects the target column from an auto-populated dropdown, and clicks ``Run Analysis.'' The system handles the rest. Fig.~\ref{fig:interface} shows the upload page with the Titanic dataset loaded, showing the auto-detected dataset summary (891 rows, 12 columns, 7 numeric features, 5 non-numeric, 866 missing values, classification task).

\begin{figure}[t]
    \centering
    % Paste your Interface Screenshot image here
    \framebox{\parbox{0.9\columnwidth}{\centering\vspace{1.5cm}\textit{[Fig. 7 -- Web Interface Upload Screen]}\vspace{1.5cm}}}
    \caption{Web interface upload page with Titanic dataset loaded. The system auto-detects 891 rows, 12 columns, and classification task type from the target column cardinality before any preprocessing runs.}
    \label{fig:interface}
\end{figure}

The preprocessing log, shown in Fig.~\ref{fig:preplog}, lists every decision the pipeline makes, in order. The log for the Titanic run shows: loaded 891 rows with 12 columns; renamed columns for consistency; dropped the ID column \texttt{passengerid}; dropped high-cardinality text column \texttt{name}; label-encoded four categorical columns (sex, ticket, cabin, embarked); imputed 177 missing values with median and mode; removed 16 duplicate rows; auto-detected classification task with target \texttt{survived} having 2 unique values; final dataset 875 rows by 10 columns. This log makes the connection between data quality decisions and the feature set that models see explicit. Students examining the log can trace each step and understand why the feature count changed from 12 to 9.

\begin{figure}[t]
    \centering
    % Paste your Preprocessing Log image here
    \framebox{\parbox{0.9\columnwidth}{\centering\vspace{1.5cm}\textit{[Fig. 8 -- Preprocessing Log]}\vspace{1.5cm}}}
    \caption{Preprocessing log for the Titanic dataset. Each step is logged with its outcome. The 16 duplicate rows removed and 177 missing values imputed are visible without any user configuration.}
    \label{fig:preplog}
\end{figure}

The recommendation panel (Fig.~\ref{fig:recommend}) identifies the top-ranked model (SVM\_rbf), reports its test accuracy (0.8114) and test F1 (0.8106), lists the 6 overfit models with a plain-language explanation, and summarises the overall fit distribution. This panel is designed for users who want a single actionable answer and do not want to interpret the full chart suite themselves.

\begin{figure}[t]
    \centering
    % Paste your Recommendation Panel image here
    \framebox{\parbox{0.9\columnwidth}{\centering\vspace{1.5cm}\textit{[Fig. 9 -- Recommendation Panel]}\vspace{1.5cm}}}
    \caption{Automated recommendation panel for the Titanic experiment. SVM (RBF) is identified as the top model (F1=0.8106, good fit). Six overfit models are listed with an explanation of the high-variance diagnosis.}
    \label{fig:recommend}
\end{figure}

\subsection{Framework Generalisability}

The preprocessing pipeline handles a broad range of tabular datasets without modification, as demonstrated by consistent operation across Titanic, heart disease, and house price datasets representing different sizes, dimensionalities, class distributions, and task types. The model catalogue covers the primary algorithm families in the scikit-learn ecosystem. One design decision worth noting is the use of fixed hyperparameter configurations for most models, with only depth (decision trees) and $k$ (KNN) varied for the complexity curve. This ensures comparability within algorithm families but means that ensemble methods may not be at their optimal configuration for a given dataset.

The framework was also informally tested on a student activity dataset (450 records, 12 features, multiclass target with 5 classes) and a customer churn dataset (7,000 records, 20 features, binary target). Both datasets produced consistent and interpretable bias-variance analyses. The churn dataset, being substantially larger than the benchmark datasets, showed narrower generalisation gaps across all models, consistent with the theoretical prediction that variance decreases with training set size.

\subsection{Limitations}

Several limitations of the current system should be acknowledged. First, the evaluation uses a single train-test split. While the fixed seed ensures reproducibility, performance differences between models with similar test F1 scores may not be statistically significant. Cross-validation support is planned as a near-term extension.

Second, the preprocessing pipeline applies a fixed strategy that may not be optimal for all datasets. Datasets with ordinal categorical variables, or where missing values are not missing at random, may benefit from more sophisticated approaches such as one-hot encoding or multiple imputation. The system does not currently detect or handle class imbalance through resampling strategies such as SMOTE~\cite{b35}, which may affect the quality of the bias-variance analysis on severely imbalanced datasets.

Third, the current system does not support text, image, or time-series data. Extension to these data types would require dataset-type-specific preprocessing modules and model families beyond the current scope.

Fourth, the bias-variance decomposition for regression tasks is estimated via bootstrap resampling, which is an approximation. With $B=50$ bootstrap samples, the estimates have non-negligible variance, particularly on small training sets. Increasing $B$ improves estimate quality but increases computation time proportionally.

Fifth, the catalogue of models, while broad, does not include neural networks. Modern deep learning models for tabular data such as TabNet~\cite{b36} or FT-Transformer exhibit bias-variance profiles that differ from the classical algorithm families and are not represented in the current analysis. Including them would require integration of PyTorch or TensorFlow as additional dependencies, which was deliberately excluded to maintain the lightweight, easy-to-install character of the framework.

\subsection{Educational Value}

The system has practical pedagogical value for machine learning courses. The decision tree complexity curve and the bias-variance bar chart with explicit GOOD/OVERFIT labels provide visual anchors for concepts that students cite as difficult to internalise from theory alone. Several features of the design reinforce this.

First, the complexity curve is generated on the user's own dataset, not on a synthetic example. This matters because the depth at which overfitting begins depends on the training set size, the number of informative features, and the class separability of the specific problem. A student who sees the curve peak at depth 3 on their 300-row medical dataset, and then at depth 5 on their 700-row passenger dataset, builds an understanding of how dataset composition affects the tradeoff. That understanding is harder to build from a fixed textbook figure.

Second, the preprocessing log makes the pipeline transparent. Students can see that removing one column and imputing another changes the feature count from 12 to 9, and then trace how the feature count change affects model performance. This connection between data decisions and downstream performance is a core practitioner skill that is infrequently taught explicitly.

Third, the model recommendation combines a ranked metric table with a plain-language explanation of the winner. On the Titanic dataset, the recommendation reads: SVM (RBF kernel) achieves the highest test F1 of 0.8106 with train accuracy 0.8443 and test accuracy 0.8114, a generalisation gap of 0.033 within the GOOD FIT threshold. It classifies unseen data nearly as well as training data. Six models are flagged as high-variance and would likely perform differently on a new split of the same data. This kind of explicit reasoning is what instructors spend class time trying to teach students to produce themselves, and seeing it generated automatically from their own experiment helps them understand what the reasoning should look like.

Student feedback from a pilot deployment in an introductory machine learning course indicated that the complexity curve was the single most useful output for building intuition about overfitting. The preprocessing log was rated the most useful output for understanding why model performance differs from naive expectations.

%% ============================================================
\section{Ablation Study: Impact of Preprocessing Choices}
%% ============================================================

\subsection{Motivation and Design}

The automated preprocessing pipeline makes several design choices that could affect the quality of the bias-variance analysis. This section presents an ablation study examining the impact of three key choices: (1) whether to standardise features; (2) whether to remove ID-like columns; and (3) the imputation strategy for missing values. The ablation is conducted on the Titanic dataset to allow direct comparison with the primary results.

Each ablation removes or modifies one preprocessing step while keeping all others fixed. The modified pipeline is applied to the same 80/20 split, and the same 15-model catalogue is trained and evaluated. The results are compared to the baseline configuration from Section VII.

\subsection{Effect of Feature Standardisation}

Removing feature standardisation from the pipeline while keeping all other steps fixed has no effect on tree-based models (decision trees, random forests, gradient boosting), which are invariant to monotone transformations of the feature values. However, it has a substantial effect on distance-based and linear models. KNN at $k=1$ achieves the same perfect training accuracy (since it memorises all training points regardless of scale) but shows a larger test accuracy gap when features are unstandardised, because the distance metric is dominated by high-scale features such as fare (range 0--512) over binary features such as sex (range 0--1). SVM test accuracy drops by approximately 3 to 5 percentage points without standardisation, and Logistic Regression shows a similar degradation.

This ablation confirms that feature standardisation is not a neutral preprocessing step: it directly affects the quality of the bias-variance analysis for scale-sensitive models. The system's default of always applying standardisation is therefore a substantive design choice with measurable consequences.

\subsection{Effect of ID-Column Removal}

If the PassengerId column (a unique integer identifier with no predictive value) is retained in the feature set, tree-based models that do not limit depth will split on PassengerId to achieve perfect training accuracy. The unconstrained decision tree achieves a training accuracy of 1.00 regardless of whether PassengerId is included, but with it included, the test accuracy drops by approximately 4 percentage points because the tree wastes splitting capacity on the meaningless identifier. The generalisation gap increases accordingly.

KNN at $k=1$ also deteriorates without ID column removal, because PassengerId is approximately linearly related to passenger boarding order and creates spurious nearest-neighbour relationships between passengers with similar IDs but different actual characteristics. This ablation confirms that ID-column removal is a substantive design choice rather than merely a cosmetic cleaning step.

\subsection{Effect of Imputation Strategy}

Three imputation strategies are compared for the missing age values: median imputation (the system default), mean imputation, and constant imputation (filling missing values with -1). For continuous features with approximately symmetric distributions, median and mean imputation produce nearly identical results. For skewed distributions, median imputation is more robust. Constant imputation with -1 produces a measurable degradation for all model families, because the constant -1 is far outside the natural age range (0--80) and creates a spurious cluster in feature space. KNN is most affected, with test accuracy dropping by approximately 6 percentage points under constant imputation.

The results confirm that median imputation is the appropriate default for continuous features and justify the system's design choice. They also demonstrate that imputation strategy is not a minor detail: it has measurable consequences on model performance and bias-variance behaviour.

\subsection{Effect of Neighbourhood Size ($k$) on KNN Variance}

The KNN complexity curve, which plots training and test accuracy as a function of $k$, shows a pattern that mirrors the decision tree depth curve but runs in reverse: decreasing $k$ from 10 to 1 increases complexity rather than decreasing it. Fig.~\ref{fig:bv_chart} from the Titanic experiment shows that KNN at $k=1$ achieves training accuracy of 0.9986 (near-perfect memorisation) but test accuracy of 0.7543, a gap of 0.244. Increasing $k$ to 10 closes this gap substantially: train accuracy falls to 0.8414 and test accuracy rises to 0.7886, a gap of 0.053, still just above the GOOD FIT threshold.

On the wine quality dataset, the pattern is more pronounced. KNN at $k=1$ achieves training accuracy 1.0000 and test accuracy 0.5245 (gap: 0.476). KNN at $k=10$ achieves training accuracy 0.6388 and test accuracy 0.5637 (gap: 0.080, classified GOOD FIT). The order-of-magnitude reduction in generalisation gap from $k=1$ to $k=10$ on a six-class dataset illustrates how a single integer hyperparameter determines whether a model overfits or generalises. The KNN complexity curve in Fig.~\ref{fig:bv_chart} makes this monotonic relationship visible by plotting the training-test accuracy gap as a function of $k$, presented left to right from high $k$ (low complexity, low variance) to low $k$ (high complexity, high variance).

Removing feature standardisation from the pipeline before training KNN inflates the gap further at every $k$ value because the distance metric is then dominated by high-range features. This interaction between preprocessing and KNN generalisation gap is one of the reasons the ablation study examines both steps together.

\subsection{Summary of Ablation Results}

Table~\ref{tab:ablation} summarises the effect of each ablation on the number of models classified as OVERFIT under each modified pipeline. The baseline configuration produces 6 OVERFIT models. Removing standardisation increases this to 9 (due to KNN and SVM degradation). Retaining the ID column increases it to 8 (due to spurious splits in tree models). Using constant imputation increases it to 10. These results confirm that each of the three preprocessing choices under study contributes meaningfully to the quality of the analysis.

\begin{table}[t]
    \caption{Ablation Study: Models Classified as OVERFIT Under Each Preprocessing Variant}
    \label{tab:ablation}
    \centering
    \small
    \renewcommand{\arraystretch}{1.3}
    \begin{tabular}{|>{\raggedright\arraybackslash}p{1.6in}|c|c|}
    \hline
    \textbf{Preprocessing Variant} & \textbf{\# OVERFIT} & \textbf{Change} \\
    \hline
    Baseline (full pipeline)        & 6  & -- \\
    \hline
    No feature standardisation      & 9  & +3 \\
    \hline
    ID column retained              & 8  & +2 \\
    \hline
    Constant imputation ($-1$)      & 10 & +4 \\
    \hline
    Mean imputation (vs.\ median)   & 6  & 0 \\
    \hline
    No duplicate removal            & 6  & 0 \\
    \hline
    \end{tabular}
\end{table}

%% ============================================================
\section{Educational Modality and Pedagogical Outcomes}
%% ============================================================

\subsection{Cognitive Scaffolding in Machine Learning Education}

The transition from theoretical machine learning constructs to applied model deployment consistently presents a steep cognitive gradient for novice practitioners. Traditional curricula frequently compartmentalise mathematical formalisms, algorithmic implementations, and empirical evaluation into distinct, poorly integrated silos. This architectural separation often forces students to manually bridge the conceptual gap between an abstract loss function differential and a tangible confusion matrix output.

The present framework systematically dismantles this barrier by providing dynamic cognitive scaffolding. By tightly coupling rigorous hyperparameter modification with instantaneous, automated graphical feedback, the system enforces an empirical feedback loop. When a student drastically decreases $k$ in a K-Nearest Neighbours classifier and immediately observes the chaotic, fragmented boundary topography and the disproportionate spike in testing error against the training error, the abstract concept of high-variance overfitting physically crystallises into a tangible empirical observation.

\subsection{Democratisation of Experimental Methodology}

Historically, robust algorithm benchmarking required extensive boilerplate coding traversing data cleaning, stratified splitting, scaling, inference, and metric aggregation. This steep engineering overhead frequently discouraged exhaustive algorithmic profiling, pushing students toward singular, familiar models like Logistic Regression or default Random Forests regardless of the underlying structural matrix topology.

The framework drastically democratises this experimental methodology. By completely automating the rigorous standardisation and preprocessing pipelines without obfuscating the underlying metrics, users can physically test fifteen independent model architectures across severely diverse geometric topographies (sparse classification, dense classification, continuous regression) within seconds. This frictionless experimentation directly encourages extensive hypothesis formulation: users can actively test whether an SVM with an RBF kernel mathematically outperforms a Gradient Boosting machine on the dense multi-class Wine Quality manifold, bypassing hours of procedural pipeline engineering.

%% ============================================================
\section{Conclusion and Future Work}
%% ============================================================

This paper presented the ML Bias-Variance Tradeoff Analyzer, an automated web-based tool for systematic evaluation of the bias-variance tradeoff across multiple machine learning classifiers and regressors on any tabular dataset. The system accepts any tabular CSV, applies a fully automated preprocessing pipeline with strict data leakage prevention, trains 15 to 18 models across six algorithm families, and produces a suite of visualisations and ranked metric tables that make overfitting, good fit, and underfitting behaviour directly visible.

For classification tasks, the system produces: a bias-variance grouped bar chart comparing train and test accuracy across all classifiers; an F1 score comparison chart; a test accuracy ranking chart; a decision tree depth-vs-accuracy complexity curve; and an SVM confusion matrix. For regression tasks, it produces a comprehensive metric table and a bootstrap-estimated bias-variance decomposition chart covering all regression models in the spectrum.

Results across four benchmark datasets (Titanic, heart disease, house price, and wine quality) confirm expected theoretical patterns and demonstrate consistent behaviour. Unconstrained models memorise training data while failing on test sets; shallow models underfit; well-regularised models at intermediate complexity achieve the best generalisation. SVM with the RBF kernel achieves the strongest overall classification performance across all three classification benchmarks, with generalisation gaps of 0.033, 0.043, and 0.085 on Titanic, heart disease, and wine quality, respectively. Gradient Boosting Regressor achieves the strongest regression performance on the Ames Housing dataset at test $R^2 = 0.887$. A shallow decision tree matches ensemble test accuracy on binary classification with a smaller generalisation gap and full interpretability, reinforcing the practical value of the complexity curve as a model selection aid.

The ablation study shows that preprocessing choices are substantive: removing feature standardisation, retaining ID columns, or using constant imputation measurably increases the number of models classified as overfit, producing a misleading picture of the bias-variance landscape. Mean imputation and duplicate removal, by contrast, have negligible effect on the Titanic dataset, showing that not all preprocessing decisions carry equal weight.

Future work will extend the framework in several directions. Cross-validation support will be integrated to provide statistically rigorous metric estimates with confidence intervals~\cite{b39}. Hyperparameter optimisation will be added to allow fair comparisons at each algorithm family's optimal configuration. The bootstrap bias-variance decomposition will be refined and extended to classification tasks following the approach of Kohavi and Wolpert~\cite{b40}. A benchmark evaluation across a broader suite of UCI and OpenML datasets~\cite{b41} will systematically characterise how dataset properties (size, dimensionality, class imbalance, feature type distribution) affect the bias-variance profiles of each model family. Integration with explainability tools such as SHAP and LIME~\cite{b42} will enable feature-level interpretation of model predictions within the same unified interface. Support for class-imbalance handling through SMOTE and class-weighted training will be added to improve analysis quality on heavily imbalanced datasets.

%% ============================================================
\begin{thebibliography}{00}

\bibitem{b1} S. Geman, E. Bienenstock, and R. Doursat, ``Neural networks and the bias/variance dilemma,'' \textit{Neural Computation}, vol. 4, no. 1, pp. 1--58, Jan. 1992, doi: 10.1162/neco.1992.4.1.1.

\bibitem{b2} J. H. Friedman, ``On bias, variance, 0/1-loss, and the curse-of-dimensionality,'' \textit{Data Mining and Knowledge Discovery}, vol. 1, no. 1, pp. 55--77, 1997, doi: 10.1023/A:1009778005914.

\bibitem{b3} F. Pedregosa et al., ``Scikit-learn: Machine Learning in Python,'' \textit{Journal of Machine Learning Research}, vol. 12, pp. 2825--2830, 2011.

\bibitem{b4} C. Thornton, F. Hutter, H. H. Hoos, and K. Leyton-Brown, ``Auto-WEKA: Combined Selection and Hyperparameter Optimization of Classification Algorithms,'' in \textit{Proc. 19th ACM SIGKDD Int. Conf. Knowledge Discovery and Data Mining}, Chicago, IL, 2013, pp. 847--855.

\bibitem{b5} M. Feurer, A. Klein, K. Eggensperger, J. Springenberg, M. Blum, and F. Hutter, ``Efficient and Robust Automated Machine Learning,'' in \textit{Advances in Neural Information Processing Systems}, vol. 28, 2015, pp. 2962--2970.

\bibitem{b6} A. Singh, S. Saraswat, and N. Faujdar, ``Analyzing Titanic disaster using machine learning algorithms,'' in \textit{Proc. 2017 Int. Conf. Computing, Communication and Automation (ICCCA)}, Greater Noida, India, 2017, pp. 406--411, doi: 10.1109/CCAA.2017.8229874.

\bibitem{b7} A. Dasgupta, V. P. Mishra, S. Jha, B. Singh, and V. K. Shukla, ``Predicting the Likelihood of Survival of Titanic's Passengers by Machine Learning,'' in \textit{Proc. 2021 IEEE Int. Conf. Computational Intelligence and Knowledge Economy (ICCIKE)}, Dubai, UAE, 2021, pp. 52--57, doi: 10.1109/ICCIKE51210.2021.9410757.

\bibitem{b8} S. Shekhar, D. Arora, and P. Sharma, ``Classifying Titanic Passenger Data and Prediction of Survival from Disaster,'' in \textit{Advances in Information Communication Technology and Computing}, Lecture Notes in Networks and Systems, vol. 135, Springer, Singapore, 2021, pp. 181--187.

\bibitem{b9} J. R. Quinlan, ``Induction of decision trees,'' \textit{Machine Learning}, vol. 1, no. 1, pp. 81--106, 1986, doi: 10.1007/BF00116251.

\bibitem{b10} C. Cortes and V. Vapnik, ``Support-vector networks,'' \textit{Machine Learning}, vol. 20, no. 3, pp. 273--297, 1995, doi: 10.1007/BF00994018.

\bibitem{b11} T. M. Cover and P. E. Hart, ``Nearest neighbor pattern classification,'' \textit{IEEE Transactions on Information Theory}, vol. 13, no. 1, pp. 21--27, Jan. 1967, doi: 10.1109/TIT.1967.1053964.

\bibitem{b12} L. Breiman, ``Random forests,'' \textit{Machine Learning}, vol. 45, no. 1, pp. 5--32, 2001, doi: 10.1023/A:1010933404324.

\bibitem{b13} J. H. Friedman, ``Greedy function approximation: A gradient boosting machine,'' \textit{The Annals of Statistics}, vol. 29, no. 5, pp. 1189--1232, Oct. 2001, doi: 10.1214/aos/1013203451.

\bibitem{b14} T. Hastie, R. Tibshirani, and J. Friedman, \textit{The Elements of Statistical Learning: Data Mining, Inference, and Prediction}, 2nd ed. New York: Springer, 2009.

\bibitem{b15} A. Tabbakh, J. K. Rout, and M. Rout, ``Analysis and Prediction of the Survival of Titanic Passengers Using Machine Learning,'' in \textit{Proc. Int. Conf. Computing and Network Communications}, Advances in Intelligent Systems and Computing, vol. 1276, Springer, Singapore, 2021, pp. 297--304.

\bibitem{b16} H. Ranglani, ``Empirical Analysis of the Bias-Variance Tradeoff Across Machine Learning Models,'' Social Science Research Network (SSRN), Preprint, Feb. 2025, doi: 10.2139/ssrn.5086450.

\bibitem{b17} K. Singh, R. Nagpal, and R. Sehgal, ``Exploratory data analysis and machine learning on Titanic disaster dataset,'' in \textit{Proc. 2020 10th Int. Conf. Cloud Computing, Data Science and Engineering (Confluence)}, Noida, India, 2020, pp. 796--801.

\bibitem{b18} L. Breiman, ``Bagging predictors,'' \textit{Machine Learning}, vol. 24, no. 2, pp. 123--140, 1996, doi: 10.1007/BF00058655.

\bibitem{b19} P. Domingos, ``A few useful things to know about machine learning,'' \textit{Communications of the ACM}, vol. 55, no. 10, pp. 78--87, Oct. 2012, doi: 10.1145/2347736.2347755.

\bibitem{b20} D. Aha, D. Kibler, and M. K. Albert, ``Instance-based learning algorithms,'' \textit{Machine Learning}, vol. 6, no. 1, pp. 37--66, Jan. 1991, doi: 10.1007/BF00153759.

\bibitem{b21} B. Belkin, D. Hsu, S. Ma, and S. Mandal, ``Reconciling modern machine-learning practice and the classical bias-variance trade-off,'' \textit{Proc. Natl. Acad. Sci. U.S.A.}, vol. 116, no. 32, pp. 15849--15854, Jul. 2019, doi: 10.1073/pnas.1903070116.

\bibitem{b22} S. Raschka, ``Model evaluation, model selection, and algorithm selection in machine learning,'' arXiv:1811.12808, Nov. 2018. [Online]. Available: https://arxiv.org/abs/1811.12808

\bibitem{b23} B. Neal et al., ``A modern take on the bias-variance tradeoff in neural networks,'' arXiv:1810.08591, Oct. 2018. [Online]. Available: https://arxiv.org/abs/1810.08591

\bibitem{b24} J. Cervantes, F. Garcia-Lamont, L. Rodriguez-Mazahua, and A. Lopez, ``A comprehensive survey on support vector machine classification: Applications, challenges and trends,'' \textit{Neurocomputing}, vol. 408, pp. 189--215, Sep. 2020, doi: 10.1016/j.neucom.2019.10.118.

\bibitem{b25} T. Chen and C. Guestrin, ``XGBoost: A scalable tree boosting system,'' in \textit{Proc. 22nd ACM SIGKDD Int. Conf. Knowledge Discovery and Data Mining}, San Francisco, CA, USA, Aug. 2016, pp. 785--794, doi: 10.1145/2939672.2939785.

\bibitem{b26} Y. LeCun, Y. Bengio, and G. Hinton, ``Deep learning,'' \textit{Nature}, vol. 521, pp. 436--444, May 2015, doi: 10.1038/nature14539.

\bibitem{b27} B. Efron and R. J. Tibshirani, \textit{An Introduction to the Bootstrap}. New York: Chapman and Hall/CRC, 1993.

\bibitem{b28} R. S. Olson and J. H. Moore, ``TPOT: A tree-based pipeline optimization tool for automating machine learning,'' in \textit{Proc. Workshop on Automatic Machine Learning, ICML}, 2016, pp. 66--74.

\bibitem{b29} E. LeDell and S. Poirier, ``H2O AutoML: Scalable Automatic Machine Learning,'' in \textit{Proc. 7th ICML Workshop on Automated Machine Learning}, 2020.

\bibitem{b30} M. Hall, E. Frank, G. Holmes, B. Pfahringer, P. Reutemann, and I. H. Witten, ``The WEKA data mining software: An update,'' \textit{ACM SIGKDD Explorations Newsletter}, vol. 11, no. 1, pp. 10--18, 2009, doi: 10.1145/1656274.1656278.

\bibitem{b31} J. Demsar et al., ``Orange: Data Mining Toolbox in Python,'' \textit{Journal of Machine Learning Research}, vol. 14, pp. 2349--2353, 2013.

\bibitem{b32} P. Flach, \textit{Machine Learning: The Art and Science of Algorithms That Make Sense of Data}. Cambridge: Cambridge University Press, 2012.

\bibitem{b33} R. Detrano et al., ``International application of a new probability algorithm for the diagnosis of coronary artery disease,'' \textit{The American Journal of Cardiology}, vol. 64, no. 5, pp. 304--310, Aug. 1989, doi: 10.1016/0002-9149(89)90524-9.

\bibitem{b34} D. De Cock, ``Ames, Iowa: Alternative to the Boston housing data as an end of semester regression project,'' \textit{Journal of Statistics Education}, vol. 19, no. 3, 2011, doi: 10.1080/10691898.2011.11889627.

\bibitem{b35} N. V. Chawla, K. W. Bowyer, L. O. Hall, and W. P. Kegelmeyer, ``SMOTE: Synthetic minority over-sampling technique,'' \textit{Journal of Artificial Intelligence Research}, vol. 16, pp. 321--357, 2002, doi: 10.1613/jair.953.

\bibitem{b36} S. O. Arik and T. Pfister, ``TabNet: Attentive interpretable tabular learning,'' in \textit{Proc. 35th AAAI Conf. Artificial Intelligence}, 2021, pp. 6679--6687.

\bibitem{b37} P. Cortez, A. Cerdeira, F. Almeida, T. Matos, and J. Reis, ``Modeling wine preferences by data mining from physicochemical properties,'' \textit{Decision Support Systems}, vol. 47, no. 4, pp. 547--553, Nov. 2009, doi: 10.1016/j.dss.2009.05.016.

\bibitem{b38} N. Gupta, A. Bhatt, P. Bansal, and S. Sharma, ``Prediction of wine quality using machine learning algorithms,'' in \textit{Proc. 2020 8th Int. Conf. Reliability, Infocom Technologies and Optimization (ICRITO)}, Noida, India, 2020, pp. 359--364, doi: 10.1109/ICRITO48877.2020.9197824.

\bibitem{b39} R. Kohavi, ``A study of cross-validation and bootstrap for accuracy estimation and model selection,'' in \textit{Proc. 14th Int. Joint Conf. Artificial Intelligence (IJCAI)}, Montreal, Canada, 1995, vol. 2, pp. 1137--1143.

\bibitem{b40} R. Kohavi and D. H. Wolpert, ``Bias plus variance decomposition for zero-one loss functions,'' in \textit{Proc. 13th Int. Conf. Machine Learning (ICML)}, Bari, Italy, 1996, pp. 275--283.

\bibitem{b41} J. Vanschoren, J. N. van Rijn, B. Bischl, and L. Torgo, ``OpenML: Networked science in machine learning,'' \textit{ACM SIGKDD Explorations Newsletter}, vol. 15, no. 2, pp. 49--60, Jun. 2014, doi: 10.1145/2641190.2641198.

\bibitem{b42} C. Molnar, \textit{Interpretable Machine Learning: A Guide for Making Black Box Models Explainable}, 2nd ed. Munich: Independently published, 2022. [Online]. Available: https://christophm.github.io/interpretable-ml-book/

\end{thebibliography}

\end{document}
